% BirdCLEF+ 2026 — Working Notes
% Last updated: 2026-04-27
% Production: v33 LB 0.932
% Total LB submissions: 21 over ~2 months
%
% =====================================================================
% EXECUTIVE SUMMARY
% =====================================================================
%
% This document is the running working-note record for our BirdCLEF+
% 2026 submission. It accumulates ~95 experiments across data pipeline,
% SED training, ensemble blending, latent diagnostics, and 21 LB
% submissions.
%
% Final state: production v33 LB 0.932, ranked ~200th out of competition
% leaderboard (top score ~0.95). Our system is Perch v2 + exp50 SED + V9
% taxon gate + file-max temporal coherence + Gaussian smoothing.
%
% Of 21 LB-submitted modifications, only 4 were positive (+0.001 each):
% v23/v24 (SED swap to exp50), v26 (weight tuning), v33 (file-max).
% Every other modification regressed by −0.002 to −0.022.
%
% Key findings driving the negative-result thesis of this work:
%
% 1. SITE-INVARIANCE HIERARCHY DETERMINES LB TRANSFER. Perch's xeno-
%    canto multi-region pretraining is the only genuinely site-
%    invariant component in our pipeline. Any modification that
%    reduces W_Perch below 0.7, or that amplifies SED-derived
%    (site-conflated) signals, regresses on LB.
%
% 2. LOCAL-LB ANTI-CORRELATION STRENGTHENS WITH MODIFICATION
%    SOPHISTICATION. Simple universal-physics levers (file-max
%    coherence at α=0.10) transferred. Per-class fitted levers (LR
%    detectors with LOSO-site CV AUC 0.97) categorically did not.
%
% 3. SITE SHORTCUT IS IN THE LABELS, NOT THE MODEL. Insecta sonotypes
%    are entirely site-conflated in our 4-Insecta-site labeled SS pool;
%    no architecture or representation can recover site-invariant
%    Insecta detection from this label structure.
%
% 4. PERCH'S LATENT KNOWS WHAT IT DOESN'T KNOW BUT THE OUTPUT HEAD
%    SUPPRESSES THIS. UNK rows show reduced embedding magnitude and
%    distinct latent geometry, but output entropy is unchanged
%    (often actively overconfident on a wrong-vocabulary species).
%
% 5. ARCHITECTURE DIVERSITY IS LIMITED BY RECIPE CONVERGENCE. Our 5
%    SED variants achieve Pearson 0.97--0.99 with each other.
%    Recipe/data tweaks on the same Boredom HGNet template produce
%    1--3% residual diversity. To break v33's ceiling, fundamentally
%    different teacher families are required (Perch v3 release wait,
%    or training from scratch on different architecture).
%
% =====================================================================
% TABLE OF CONTENTS (sections numbered chronologically as added)
% =====================================================================
%
%   §1  Background and v12 baseline (LB 0.929)
%   §2  Pipeline failure-mode audit (exp45)
%   §3  exp47--50 Boredom recipe + 2025 BG → v24/v26 (+0.001 each LB)
%   §4  exp51 27-class head + v28 LB regression
%   §5  exp52 v22 catastrophe post-mortem
%   §6  exp53--54 lever exploration (53 family negative, 54 positive)
%   §7  exp55--58 mechanism analyses
%   §8  exp59--69 architecture diversity attempts
%   §9  exp70--75 audit cycle
%   §10 exp76--80 iVAE / mel-iVAE chain (v34 LB regression, exp80 teardown)
%   §11 exp81--85 Q1--Q4 systematic local audit
%   §12 v36 5-way blend LB regression (exp82c-85)
%   §13 v33 file-max temporal coherence (LB +0.001)
%   §14 0.929 → 0.932 honest decomposition (3 site-invariance sources)
%   §15 exp89--94 latent FFT + TP/FP/TN/FN signal hunt
%   §16 exp95--98 multi-feature signal mining + S08 deep dive
%   §17 exp99--100 LR-correction lever (LOSO-site CV 0.97)
%   §18 exp101--102 multi-teacher variants + LR artifacts
%   §19 Final 5-slot LB sweep (2026-04-27, all 5 regressed)
%   §20 Negative-result thesis: why all novel attempts failed
%
% =====================================================================
% FINAL LB SUBMISSION TIMELINE
% =====================================================================
%
%   v1   Perch only (frozen probe, no SED)              0.910
%   v6   0.80 P + 0.20 SED29 + Gauss σ=0.5              0.912/0.927
%   v12  + V9 taxon gate                                 0.929
%   v16  SED41f alone (no Perch)                         0.856
%   v17  0.8P + 0.2 SED41f                               0.922
%   v18  v17 + exp44c overlay                            0.916
%   v19  v17 + exp44g synth                              0.907
%   v20  v17 + V9 gate                                   0.929
%   v21  v12 + V9 gate                                   0.929
%   v22  + exp48 site prior + cluster                    0.913
%   v23  0.8P + 0.1 SED29 + 0.1 exp50                   0.928
%   v24  0.8P + 0.2 exp50                                0.930
%   v25  0.5P + 0.5 exp50                                0.928
%   v26  0.7P + 0.3 exp50  (peak SED weight)             0.931
%   v28  v26 + exp51 27-head additive                    0.929
%   v29  v26 + per-class OOF routing                     0.927
%   v32  v26 + R5 aliveness routing                      0.930
%   v33  v26 + L2c file-max coherence α=0.10  PRODUCTION 0.932
%   v34  v33 + mel-iVAE z-kNN                            0.916
%   v36  5-way blend W_Perch=0.5                         0.915
%   v38  LR α=0.3 β=0.1                                  0.920
%   v39  LR α=0.5 β=0.1                                  0.915
%   v40  LR α=0.5 β=0.0 (FN-only)                        0.917
%   v41  threshold rule (no LR)                          0.918
%   v42  file-max α=0.20                                 0.918
%
% =====================================================================

% exp_current.tex — appendix of exp37-43 findings for working-note draft.
% Merge into experiments.tex when stable. Do not edit experiments.tex directly
% for in-flight exp42/43 work; stage here to keep main narrative clean.

\subsection{Experiments 37--40: SED retraining under fair evaluation}
\label{sec:exp37-40}

The 2025-winner recipe (Salman, 18th, LB 0.922 single-SED) lists five key
ingredients: HGNetV2-B0 backbone, 20-second chunks, raw-waveform mixup,
cross-entropy on clipwise + framewise-max logits, and 20+ epochs.
We faithfully re-executed a sweep across these ingredients and recorded that
no combination beat our SED29 baseline alone on the \emph{fair} 11-file
held-out evaluation (Val-A\textsubscript{v2}, 132 segments, introduced in
exp38 after identifying a 48-file training leak in the original Val-A).

\paragraph{Key negative findings.}
\begin{itemize}
  \item \textbf{exp37b:} HGNetV2-B2 backbone swap (all other exp29 conditions
    identical) collapsed from Val-A 0.737 (B0) to 0.58 (B2). The larger
    backbone overfits our 35k-sample train\_audio set.
  \item \textbf{exp37c:} Data expansion (labeled SS + 2025-overlap species)
    without architecture changes plateaued at 0.63.
  \item \textbf{exp40:} ConvNeXt-small SED ramped slowly (ep4 0.63) and
    never matched HGNet-B0; blend contribution negligible.
\end{itemize}

\paragraph{exp38 --- labeled soundscape inclusion.}
The first clear advance came from including labeled soundscape chunks in
training. With a 55/11 file split, the student (HGNet B0, 20-sec, BCE+BCE,
WeightedSampler 20\% SS ratio) reached Val-A\textsubscript{v2} 0.810 at
epoch 9. Fair-eval three-way blend \{Perch 0.70, SED29 0.15, SED38 0.15\}
achieved 0.902, +0.011 over Perch alone.

\paragraph{exp39 --- backbone diversity.}
An EfficientNet-B0 (NS-JFT pretrained) SED trained identically to exp29
achieved Val-A 0.789 alone (+0.052 over SED29). Pearson correlation with
SED29 was 0.43 (moderate diversity). Three-way blend \{Perch 0.70, SED29
0.15, SED39 0.15\} reached 0.909 on fair eval.

\paragraph{LB reality check.}
Despite local improvements, every SED-swap or SED-addition variant we
submitted to the leaderboard \emph{regressed} vs. our v12 reference
(perch-distill with SED29 blend, LB 0.929). Specifically:

\begin{center}
\begin{tabular}{llrr}
\toprule
Version & Configuration & Fair Val-A & LB \\
\midrule
v12 (ref) & Perch + SED29 (0.80) + Gauss 0.5                & 0.897 & \textbf{0.929} \\
v14       & SED29 $\to$ SED38 swap                           & 0.902 & 0.926 \\
v13       & 3-way + SED38                                    & 0.973\textsuperscript{*} & 0.925 \\
v15       & 3-way + SED39                                    & 0.909 & 0.922 \\
v17       & Perch + SED41f (ensemble-pseudo, $\S$\ref{sec:exp41}) & 0.922 & 0.922 \\
\bottomrule
\end{tabular}
\end{center}
\textsuperscript{*} v13 local Val-A contains leakage with SED38 training files.

The consistent pattern --- local positive, LB flat or negative ---
indicates the fair 11-file evaluation has high variance and is not a
reliable LB predictor. We treat it as a necessary but insufficient gate
and maintain LB submission discipline (5/day hard limit).

\subsection{Experiment 41: Pseudo-labeling pipeline}
\label{sec:exp41}

We implemented a 2025-winner-style self-training pipeline on the 10{,}592
unlabeled soundscape files. Teacher: SED29 (train\_audio only);
Student: HGNet B0 retrained from scratch.

\paragraph{Filter chain.}
Following the 1st/2nd prize write-ups of BirdCLEF 2025:
\begin{enumerate}
  \item Discard chunks where $\max_c p_c < 0.5$ (exp41c, 2nd prize recipe).
  \item Zero out class probabilities below 0.1 (1st prize recipe).
  \item Power transform $p^{0.5}$ to reduce pseudo-label noise.
  \item Per-chunk Mahalanobis confidence weight (exp41b, using SED29
    pre-head features as proxy since Perch was not yet extracted on the
    unlabeled soundscape corpus).
\end{enumerate}

\paragraph{Ensemble-teacher student (exp41f).}
The strongest variant used an ensemble of (Perch probes $\oplus$ SED29)
as teacher for pseudo generation. Student trained on train\_audio
(35{,}549) + labeled SS (55 files) + pseudo (18{,}436 chunks). Best epoch
8: Val-A\textsubscript{v2} 0.902 alone; three-way blend with Perch at
$\alpha=0.5$: \textbf{0.922} fair. This is our current local ceiling.

\paragraph{r2 regression.}
Treating exp41f predictions as a new teacher and running a second round
of pseudo-generation (exp41h) reduced Val-A\textsubscript{v2} to 0.881
alone, 0.914 blended. r2 requires a diversified teacher family
(heterogeneous backbones) to avoid confirmation collapse --- single-family
iteration runs out of signal quickly.

\subsection{Experiment 42--43: Identifiable latent pseudo-refinement}
\label{sec:exp42-43}

The core question for our working-note: can a structurally-motivated
latent space (iVAE with auxiliary conditioning) improve pseudo-label
confidence beyond the existing Mahalanobis / teacher-posterior signals?

\subsubsection{exp42 --- kNN label agreement under linear corrections}
On the 708 labeled soundscape windows with cached Perch embeddings
(exp21), cross-site kNN label agreement was measured. Raw Perch
(cosine, $k=5$, same-file excluded) achieved cross-site agreement
0.638 (vs.\ random 0.456, $+0.18$ lift) --- evidence of real species
signal orthogonal to site co-occurrence. Site-mean centering and
file-mean centering \emph{dropped} agreement to 0.513 and 0.425
respectively, consistent with the exp26 Val-B collapse finding:
linear corrections strip species signal along with the site shortcut.

\subsubsection{exp43b --- iVAE on Perch embeddings}
We implemented iVAE \citep{khemakhem2020ivae} with $u = $ one-hot(site,
hour), $\dim(z) = 32$. Our first training run suffered classical
posterior collapse (KL $\to 0$ by epoch 30, $\sigma_z \to 0.1$) due to
the standard reduction-\texttt{mean} MSE hiding a $1536\times$ scale
mismatch with the KL term. After fixing to reduction-\texttt{sum} per
window, adding $\beta$ warmup $[0\to 1]$ over 50 epochs, standardizing
Perch input, and applying free-bits (floor of 0.5 nats per dim), the
model converged cleanly to $\sigma_z \approx 2.4$, rec 226 (standardized).
Probe AUC: z-32 = 0.934 vs.\ PCA-32 = 0.935 --- iVAE z is
\emph{identity-preserving} in species discrimination.

\subsubsection{exp43c --- mislabel detection}
We injected 15\% synthetic label flips and measured ROC-AUC of each
candidate \emph{disagreement score} for ranking flipped vs.\ clean windows.
Results at $k=5$, cross-site neighbor restriction:

\begin{center}
\begin{tabular}{lr}
\toprule
Space & Flip-detection AUC \\
\midrule
Perch raw (1536-d, standardized) & 0.605 \\
PCA-32                           & 0.597 \\
iVAE z-32                        & \textbf{0.618} \\
iVAE z-32 $\oplus$ PCA-32 (64-d) & 0.605 \\
\bottomrule
\end{tabular}
\end{center}

iVAE z-32 attains the highest single-space signal at $k=5$, though the
gap over raw Perch is small ($+0.013$) and Spearman correlation with raw
Perch kNN signal is 0.84 at this scale, suggesting iVAE at
$n=708$ is largely a noisy replica of raw-Perch-kNN rather than a
structurally distinct signal.

\subsubsection{exp43d --- entropy and orthogonality analysis}
Teacher pseudo probabilities on the 10{,}592-file unlabeled corpus are
broadly uncertain: 82\% of windows have $\max p_c < 0.5$; median
$\max p_c = 0.29$. Only 5.5\% are confident ($\max p_c > 0.7$). An
ambiguous zone (top-1 minus top-2 $< 0.2$) captures 16.6\%
(21{,}140 windows) --- the natural target for a disambiguation filter.

Crucially, Spearman correlation between Mahalanobis confidence weight
(exp41b) and teacher posterior sharpness is $+0.15$ --- nearly
orthogonal. This motivates a multi-signal composite filter.

\subsubsection{exp43g --- composite confidence filter}
On labeled SS mislabel detection, we combined iVAE z-kNN disagreement,
raw-Perch kNN disagreement, teacher-ensemble posterior disagreement
($1 - P(y_\text{label} \mid x)$), and Mahalanobis distance:

\begin{center}
\begin{tabular}{lr}
\toprule
Signal / composite & Flip-AUC \\
\midrule
Teacher (1 - P) alone              & 0.687 \\
Perch-kNN alone                    & 0.606 \\
iVAE-32 kNN alone                  & 0.568 \\
Mahalanobis alone                  & 0.463 \\
Perch-kNN + Teacher (rank-mean)    & \textbf{0.707} \\
All 4 signals (rank-mean)          & 0.632 \\
\bottomrule
\end{tabular}
\end{center}

Key findings: (i) \emph{teacher posterior is the strongest single signal}
and was not being exploited by the exp41c filter chain; (ii) Mahalanobis
on this dataset performs below chance for label-error ranking and
\emph{hurts} composite scores when included; (iii) at $n=708$ the iVAE
signal is largely redundant with raw-Perch kNN (Spearman 0.84). We
expect the iVAE signal to separate from raw-Perch at scale
($n\approx 128{,}000$), which is the motivating experiment
(exp43e, in flight).

\subsubsection{exp43i --- teacher-sharpness weighting (ablation)}
We replaced the Mahalanobis weight in exp41c with a teacher-sharpness
weight $(p_\text{top1} - p_\text{top2}) / p_\text{top1}$ and retrained
the student. Fair Val-A\textsubscript{v2} 0.865 alone (vs.\ exp41f
0.898); blend 0.914 (vs.\ exp41f 0.922) --- \emph{negative}. We
interpret this as over-weighting already-confident pseudo examples and
losing useful diversity. The stronger mechanism is to use sharpness as
a filter gate rather than a raw sample weight.

\subsubsection{exp43e --- pooled iVAE at 128k fails the thesis}
We extracted Perch embeddings on the full 10{,}658-file soundscape
corpus (127{,}896 windows) via ONNX on GPU (0.81 sec/batch of 192,
$\approx 22$ min wall), then trained an MLP iVAE with $\dim(z)=32$,
$u=(\text{site}, \text{hour})$, non-zero $\beta$-warmup and free-bits.
Training converged cleanly (KL 63, $\sigma_z = 1.93$, no collapse), but
flip-detection AUC @ $k{=}10$ on labeled subset was 0.562 (iVAE) vs.
0.568 (raw Perch): \emph{zero improvement}. Spearman correlation of
kNN disagreement scores \emph{increased} from 0.84 at $n=708$ to
\textbf{0.901 at $n=128{,}000$}. Scale worsens redundancy.

Interpretation: Perch pooled embedding is already a species-optimized
compression. There is no residual site-confound left for a linear-prior
iVAE to disentangle in $\mathbb{R}^{1536}$. The paper thesis as
originally framed (iVAE on pooled features) is not viable.

\subsection{Pivot to iVDFM --- innovation-conditioned time-series iVAE}
\label{sec:ivdfm}

We therefore adopt the iVDFM framework of \citet{chang2026ivdfm}, which
adapts iVAE to multivariate time series via four changes, each of
which targets a specific failure mode of pooled-feature iVAE.

\paragraph{Innovation-level conditioning.}
Rather than conditioning the latent-state prior $p(z_t \mid u_t)$
(standard time-series iVAE extension), iVDFM conditions the
\emph{innovation} process $p(\eta_t \mid u_t, e_t)$ that drives the
dynamics. Identifiability propagates to factors through linear diagonal
dynamics without the rotational ambiguity that afflicts classical DFMs.

\paragraph{Non-Gaussian innovation prior.}
\citet{chang2026ivdfm} proves Gaussian innovations are degenerate: for
any invertible linear $R$, $\tilde \eta = R\eta$ induces the same
observation distribution. iVDFM uses Laplace (or other non-Gaussian
exponential family) innovations. In our pooled-iVAE exp43e we used
Gaussian -- this alone may explain the observed redundancy.

\paragraph{Regime embedding.}
$e_t = \sum_k \pi_{t,k} e_k$ with $\pi_t = \mathrm{softmax}(
\text{RegimeNet}(u_t))$ captures discrete acoustic regimes
(e.g.\ quiet background, bird vocalization, insect stridulation,
amphibian call) while remaining deterministic given $u_t$. $K=4$ in our
BirdCLEF adaptation.

\paragraph{Diagonal dynamics.}
$f_{t+1} = \bar A_t f_t + \bar B_t \eta_t$ with diagonal $\bar A, \bar
B$ (regime-mixed). The diagonal structure preserves the component-wise
equivalence class from innovation prior to factor trajectory.

\paragraph{BirdCLEF adaptation (exp43k).}
Input: Perch \texttt{spatial\_embedding} output, shape $(B, 16, 4, 1536)$,
averaged over the 4 frequency patches to $(B, T{=}16, 1536)$ per
5-second window (16 temporal patches of $\approx 313$ ms).
Auxiliary $u_t = (\mathrm{site\_oh}_{23} \oplus \mathrm{hour\_oh}_{18}
\oplus \mathrm{pos\_oh}_{16}) \in \mathbb{R}^{57}$; the
within-window position ensures $u_t$ varies across $t$, providing
sufficient auxiliary variability for identifiability
($\lvert (u, e)\rvert \gg r+1$).
Innovation $\eta_t \in \mathbb{R}^{16}$, Laplace.
Regime $K=4$, regime code $e_k \in \mathbb{R}^{16}$.
Encoder: 2-layer bidirectional GRU over $(y_t \oplus u_t \oplus e_t)$.
Decoder: 2-layer MLP, output $\hat y_t \in \mathbb{R}^{1536}$.
Training: 25 epochs AdamW $3\times10^{-4}$, cosine schedule, batch 128,
$\beta$-warmup 0-1 over 8 epochs, free-bits 0.3 nats.

\paragraph{Gates for paper viability.}
\begin{enumerate}
  \item[G1] Flip-detection AUC @ $k{=}10$ on labeled subset
    $\geq 0.626$ (raw-Perch-pooled baseline 0.606, +0.02 gate).
  \item[G2] Spearman correlation of iVDFM-$f$ kNN disagreement vs.
    raw-Perch-pooled kNN disagreement $\leq 0.50$ (distinct signal).
  \item[G3] Regime posterior $\pi_t$ separates qualitatively by coarse
    acoustic class (inspection).
\end{enumerate}

G1 and G2 together establish that the identifiability guarantee of
iVDFM produces a genuinely new signal for pseudo-label refinement
beyond what a pre-trained foundation model (Perch) already supplies.
G3 supports interpretability claims in the working-note narrative.

\subsection{Compute discipline}

All local training uses the NVIDIA RTX 5090 (32 GB). The Google Perch v2
TensorFlow SavedModel, however, is XLA-compiled for CPU only and refuses
to execute on GPU devices. For local extraction we use the ONNX export
(\emph{rishikeshjani/perch-onnx-for-birdclef-2026}) with
\texttt{CUDAExecutionProvider}: 0.81 sec/batch of 192 windows
(vs.\ 8 sec/batch TF-CPU), a 10$\times$ speedup that makes full-corpus
Perch extraction (10{,}658 files) a 20-minute operation rather than
three hours.

% =====================================================================
\subsection{Experiment 44: Dedicated double-blind species head}
\label{sec:exp44}

Motivated by the exp43r audit showing that 27 of 234 competition species
are in neither Perch's 14{,}795-class training set nor
\texttt{train\_audio} (25 Insecta sonotypes \texttt{47158sonXX}, 2 rare
Amphibia), we train a dedicated 27-class head on the labeled-soundscape
55-file training split.

\paragraph{exp44c --- baseline dedicated head.}
HGNetV2-B0 SED, input BatchNorm on the 128-mel channel dimension, raw
20-second clips with per-window attention pooling, AdamW
$1\times10^{-4}$, BCE-clipwise + BCE-framewise-max, 6$\times$ positive
oversampling, pure fp32 throughout (AMP fp16 overflows on RTX 5090
compute-capability 12.0a, producing NaN batches; we observed this
reproducibly and disabled AMP). Best
Val-A\textsubscript{v2} \textbf{0.848 @ epoch 4} on 15 of 27 evaluable
classes. Per-class AUC distribution is strongly bimodal: six species
at $>0.95$ (site-concentrated Insecta sonotypes), two at
$\approx 0.53$ (random).

\paragraph{exp44e --- blend evaluation.}
On 40 of 234 evaluable classes (11 held-out files, 122 windows), macro
AUC by overlay variant:

\begin{center}
\begin{tabular}{lccc}
\toprule
Variant & Macro (all 40) & Mapped (25) & Unmapped (15) \\
\midrule
Raw Perch           & 0.622 & 0.695 & \textbf{0.500 (random)} \\
Perch + Gauss        & 0.626 & 0.702 & 0.500 \\
Raw Perch + exp44c   & \textbf{0.752} & 0.695 & \textbf{0.848} \\
Perch + Gauss + exp44c & \textbf{0.757} & 0.702 & 0.848 \\
\bottomrule
\end{tabular}
\end{center}

Δ macro = $+0.131$ from the overlay, driven by the 27 unmapped columns
going from random (AUC 0.500) to 0.848 per-class average.

\paragraph{exp44g --- synthetic site augmentation.}
Since 16 of 27 species occur in a single site only, the baseline head
risks learning ``site fingerprint = sonotype''. We augment by mixing
each source call clip with a background clip drawn from a different
site's unlabeled soundscape pool: $\alpha \cdot \text{source} + (1 -
\alpha) \cdot \text{bg}$ with $\alpha \sim U(0.4, 0.8)$, applied with
probability $0.75$ during training. Best Val-A\textsubscript{v2}
\textbf{0.884} ($+0.036$ over exp44c). Bottom-10 per-class AUC
improves from $[0.53, 0.53, 0.71, \ldots]$ to
$[0.62, 0.67, 0.75, \ldots]$, suggesting the site shortcut is partly
broken.

\paragraph{exp44f --- leaderboard reality check.}
We integrate both variants as drop-in replacements for the 27 unmapped
columns in the \texttt{perch-distill} submission pipeline:

\begin{center}
\begin{tabular}{lcc}
\toprule
Version & Configuration & LB \\
\midrule
v12 (reference) & Perch + SED29 blend & \textbf{0.929} \\
v17 & Perch + SED41f blend (no overlay) & 0.922 \\
v18 (2 runs) & v17 + exp44c overlay & 0.916 / 0.914 \\
v19 & v17 + exp44g overlay & \textbf{0.907} \\
\bottomrule
\end{tabular}
\end{center}

Both overlays regress. Synthetic augmentation \emph{worsens} LB
despite improving local fair-eval. Our fair 11-file Val-A\textsubscript{v2}
is not only uninformative but \emph{anti-correlated} with LB for
single-site taxa: better local $\Rightarrow$ worse LB.

We attribute this to \textbf{distributional (covariate-shift)
overfitting rather than classical overfitting}. Training curves are
clean, model parameter count is modest, and loss converges. What the
model memorizes is the recording-equipment acoustic fingerprint that
co-varies with site in our training distribution. Synthetic mixing
diversifies \emph{backgrounds} but cannot remove the source's mic-and-
preamp response --- so exp44g learns to attend to this fingerprint more
confidently, and misfires catastrophically on hidden test sites. This
is a specific instance of the well-known failure mode where domain-
randomised data augmentation randomises the wrong channel
\citep{gulrajani2021search}.

\paragraph{Takeaway.} Dedicated heads with
overlay integration cannot work for single-site taxa without either
multi-site labeled data, explicit domain-adversarial training, or
external recordings that de-correlate site from species. None are
available under the stated competition data policy.

% =====================================================================
\subsection{Experiment 45: Pipeline failure-mode audit}
\label{sec:exp45}

The v18/v19 regressions prompted a systematic analysis of which species
our pipeline actually fails on, beyond the 27 known-unmapped double-
blind tail.

\paragraph{Setup.}
For each of the 234 competition classes with at least one positive and
one negative in the 11-file held-out eval, we compute per-class ROC-AUC
under four variants: raw Perch probabilities, SED29 alone, SED41f
alone, and a v12-analog z-blend $0.80\,\text{Perch} + 0.20\,\text{SED29}$
with $\sigma{=}0.5$ Gaussian smoothing. Forty classes qualified.

\paragraph{Model component comparison.}

\begin{center}
\begin{tabular}{lc}
\toprule
Variant & Macro AUC (40 classes) \\
\midrule
Raw Perch alone       & 0.622 \\
SED29 alone           & 0.715 \\
v12 z-blend           & 0.714 \\
\textbf{SED41f alone} & \textbf{0.878} \\
\bottomrule
\end{tabular}
\end{center}

SED41f dominates locally; the v12 blend does not recover this. This
recovers the local mystery of why the (Perch+SED41f) blend in v17 lost
to v12 on LB: our blend is on the locally-wrong side of the
local/LB divergence.

\paragraph{What predicts per-class AUC.}

Stratified by \texttt{train\_audio} clip count (species support in the
clean-recording corpus):

\begin{center}
\begin{tabular}{ccc}
\toprule
Clip count & $n$ classes & Mean AUC \\
\midrule
$0$             & 16 & 0.706 \\
$1$--$10$       & 6  & \textbf{0.491 (worst)} \\
$11$--$50$      & 9  & 0.765 \\
$51$--$200$     & 6  & 0.771 \\
$\geq 200$      & 3  & 0.939 \\
\bottomrule
\end{tabular}
\end{center}

\textbf{The one-to-ten-clip band performs \emph{worse} than zero
clips.} At zero clips, Perch's output on the corresponding label
column is structurally constant (the logit is always taken from
unmapped entries), so its per-class AUC is exactly 0.5. With 1--10
clips the pipeline learns a noisy signal that is \emph{confidently
incorrect} on held-out. This is the ``data-scarcity cliff'': enough
samples to fit, not enough to generalise.

\paragraph{Bottom-8 offenders and their confusion.}

\begin{center}
\begin{tabular}{llcll}
\toprule
Label & AUC & Taxon & $n_\text{TA}$ & Confused with (top-3) \\
\midrule
516975 (\emph{Mico melanurus}) & 0.000 & Mammalia  & 3   & litnig1, soulap1, whnjay1 (all Aves) \\
67107                          & 0.333 & Amphibia  & 0   & 22961, 22930, bcwfin2 (Amphibia cluster) \\
326272                         & 0.358 & Amphibia  & 0   & 22961, 22930, bcwfin2 \\
bafcur1                        & 0.398 & Aves      & 125 & strcuc1, sobtyr1, undtin1 (sister Aves) \\
74113                          & 0.400 & Mammalia  & 2   & compau, osprey, bbwduc (all Aves) \\
25073                          & 0.407 & Amphibia  & 0   & 22961, 22930, bcwfin2 \\
116570 (\emph{Caiman yacare})  & 0.441 & Reptilia  & 1   & flawar1, chobla1, yebela1 (all Aves) \\
47158son11                     & 0.464 & Insecta   & 0   & 22961, 22930, bcwfin2 \\
\bottomrule
\end{tabular}
\end{center}

Two structural patterns emerge:

\begin{enumerate}
  \item \textbf{Taxonomic inductive bias of Perch.} Every
    non-Aves non-Insecta offender (3 Mammalia + 1 Reptilia + 1
    Amphibia with non-zero $n_\text{TA}$) gets confidently classified
    as one of several bird species. Perch does not abstain on unfamiliar
    taxa; it maps them onto the nearest-sounding bird in its
    14{,}795-class training set. A monkey (\emph{Mico melanurus},
    \emph{Sapajus cay}) and a caiman (\emph{Caiman yacare}) are
    uniformly classified as bird species. This is a structural
    limitation of a bird-specialist foundation model used on
    multi-taxa wildlife monitoring.

  \item \textbf{Sister-species confusion within Aves.} The one Aves
    species in the bottom set (\texttt{bafcur1}) has abundant training
    data ($n_\text{TA}=125$) but is still confused with three other
    Aves. Standard cross-entropy on independent classes cannot exploit
    between-species acoustic similarity; without pairwise
    discriminators, close relatives collapse.
\end{enumerate}

Windows in our eval are overwhelmingly multi-label: 85\% contain 3 or
more simultaneously vocalising species. The per-class BCE assumption
of independence is violated by construction.

\paragraph{Next-lever recommendations.}

The audit points to two orthogonal levers not yet tried, each targeting
a distinct failure mode:

\begin{itemize}
  \item \textbf{Taxonomic hierarchical loss}: first predict
    \texttt{class\_name} (Aves / Amphibia / Insecta / Mammalia /
    Reptilia), then species conditional on taxon. Would directly
    attack the non-Aves $\to$ Aves confusion that dominates the
    bottom-8 list.
  \item \textbf{Dedicated Mammalia/Reptilia sub-head with data-scarcity
    mitigation}: the 1--10-clip band needs external recordings,
    spectral-augmentation that actually removes recording-chain
    fingerprints, or an explicit domain-adversarial loss. Naïve
    oversampling reproduces the exp44g failure mode.
\end{itemize}

The audit also clarifies the iVDFM story:
for single-site taxa, identifiable latent disentanglement of
\emph{site} as a nuisance factor is the principled solution, but the
auxiliary variable $u$ must be grounded in acoustic descriptors that
actually measure the recording chain, not just site identity (which
is perfectly entangled with species in the training data).

% =====================================================================
\subsection{Experiment 45a-f: Hierarchical classifier + per-class audit}
\label{sec:exp45af}

Motivated by exp45's finding that Perch systematically mispredicts
non-Aves taxa as Aves species, we trained a 5-way taxon classifier
(Aves / Amphibia / Insecta / Mammalia / Reptilia) on Perch 1536-d
embeddings with train\_audio (35{,}549) + labeled soundscapes (617 on
the 55-file training split) $= 36{,}166$ samples. A 2-layer MLP with
GELU nonlinearity, 256 hidden units, AdamW $1\times10^{-4}$, BCE loss.

\paragraph{exp45a --- multiplicative gating.}
Combination rule: $\text{final}[c] = \text{base}[c] \cdot
\text{taxon\_prob}[\tau(c)]$ where $\tau(c)$ is the taxon of class $c$.
Evaluated on 40 eval classes against raw Perch baseline (0.622 macro):
macro AUC rises to \textbf{0.745 ($+0.123$)}. Per-taxon: Aves $+0.186$,
Amphibia $+0.068$, Insecta $+0.181$, \textbf{Mammalia $-0.117$ (regresses)},
Reptilia $+0.280$. The Mammalia regression traces to species where the
taxon head itself is uncertain (e.g., 74113: 0.375 $\to$ 0.042); the
multiplicative gate zeroes out legitimate species signal when the taxon
detector fails.

\paragraph{exp45b --- class-balanced pos\_weight.}
Added per-taxon $\text{pos\_weight}_t = N / (5 \cdot N^+_t)$, capped at
100. Macro drops slightly to 0.737. Balancing does not fix the 74113
problem (species-level data scarcity is the root cause, not taxon-level
imbalance). \emph{Accepting that class balance is the wrong axis.}

\paragraph{exp45c --- soft gating variants.}
Ten combination rules tested. The best is V9, an additive-offset soft
gate: $\text{final}[c] = \text{base}[c] \cdot (\text{taxon\_prob}[\tau(c)]
+ 0.1)$. Macro climbs to \textbf{0.767 ($+0.145$ over Perch-alone)}, and
crucially \textbf{every taxon is positive} ($+0.186, +0.081, +0.181,
+0.122, +0.280$). The additive offset floors the gate so species with
low taxon probability are attenuated but not zeroed --- preserving
information when the taxon detector fails.

\paragraph{exp45d --- SED weight sweep.}
Grid sweep over $(w_P, w_{29}, w_{41\text{f}})$ with $\sigma{=}0.5$
Gaussian smoothing and z-score blend, on 40 eval classes:

\begin{center}
\begin{tabular}{lcc}
\toprule
Configuration                                  & Local macro & LB \\
\midrule
v12 ($0.8P + 0.2\text{S29}$)                   & 0.714       & \textbf{0.929} \\
v17 ($0.8P + 0.2\text{S41f}$)                  & 0.814       & 0.922 \\
3-way ($0.7P + 0.15\text{S29} + 0.15\text{S41f}$) & 0.816    & 0.925 \\
SED41f alone ($w_{41}=1$)                      & \textbf{0.881} & \textbf{0.856} (v16) \\
\bottomrule
\end{tabular}
\end{center}

Local macro and LB are strongly \textbf{anti-correlated} along the
$w_{41\text{f}}$ axis. This is not noise: SED41f was trained on
pseudo-labels generated by a Perch-centered ensemble, so
Perch's site-specific confidences are baked into SED41f. These
confidences help on the same-site held-out local eval but hurt on
hidden-site LB.

Per-species decomposition: SED41f perfectly resolves bafcur1 (Aves
AUC $0.278 \to 0.946$), litnig1 ($0.364 \to 1.000$), wfwduc1
($0.488 \to 0.996$) and 27 Insecta sonotypes ($0.500 \to 1.000$).
None of these local gains transfer to LB.

\paragraph{exp45e --- taxon gate $\times$ base interaction (critical).}
Apply the V9 gate to each candidate base blend:

\begin{center}
\begin{tabular}{lccc}
\toprule
Base            & Base macro & Gated macro & $\Delta$ \\
\midrule
Perch alone     & 0.622 & \textbf{0.767} & $+0.145$ \\
v12 ($0.8P + 0.2\text{S29}$) & 0.714 & 0.681 & $-0.033$ \\
v17 ($0.8P + 0.2\text{S41f}$) & 0.814 & 0.782 & $-0.032$ \\
3-way           & 0.816 & 0.799 & $-0.017$ \\
\bottomrule
\end{tabular}
\end{center}

\textbf{The gate helps only when the base has no species-level signal.}
On the Perch-alone base, many classes are stuck at 0.5 (no mapping or
no signal); any soft taxon signal differentiates them. But once a
competent SED is in the blend, the SED provides real species-level
ranking that the multiplicative taxon-gate disturbs. For bafcur1, where
SED41f reached 0.95 on its own, applying the gate dropped it to 0.14 ---
a catastrophic noise injection.

This reframes the exp45a/c headline ``+0.145 macro from taxon gating''
as an artifact of using Perch-alone (a strawman base) for comparison.
Against a production base (Perch + SED blend), the gate is a loss.

\paragraph{exp45f --- per-class optimal teacher weights.}
For each of 40 eval classes, we grid-search the optimal $\alpha$ in
$(1-\alpha)\,\text{Perch} + \alpha\,\text{SED41f}$. Distribution of
optimal $\alpha$ is bimodal: 22 classes prefer $\alpha < 0.25$, 15
prefer $\alpha > 0.75$, only 3 are in between.

Oracle per-class blend reaches macro 0.903 --- an $+0.025$ upper bound
over SED41f alone. Per-taxon oracle ceilings: Aves 0.965, Amphibia
0.895, Insecta 0.906, Mammalia 0.772 (data-scarcity-limited), Reptilia
0.818.

\paragraph{LB validation of all five directions.}
We submitted four variants spanning the local/LB trade-off to verify
the anti-correlation claim:

\begin{center}
\begin{tabular}{llcc}
\toprule
v  & Pipeline                             & Local $\Delta$ macro & LB \\
\midrule
v18 & v17 + exp44c overlay (single-site 27-class head) & $+0.131$ & $0.916$ ($-0.013$) \\
v19 & v17 + exp44g synthetic augmentation              & $+0.167$ & $0.907$ ($-0.022$) \\
v20 & v17 + V9 taxon gate                              & $-0.032$ & \textbf{$0.929$} ($+0.007$) \\
v21 & v12 + V9 taxon gate                              & $-0.033$ & \textbf{$0.929$} ($0.000$) \\
\bottomrule
\end{tabular}
\end{center}

The anti-correlation holds in four independent submissions. Local
\emph{negative} is required for LB \emph{non-negative} in this
setup. Interventions that raise local macro by targeting the rare/hard
40 eval classes must by construction disturb the 194 common classes
invisible to the local eval but dominant in the LB macro. The gate
does not push above the v12 ceiling; it can only recover drag from
suboptimal bases (v17 to v12 ceiling, $+0.007$).

\paragraph{Final lesson for this competition.}
The fair 11-file held-out is a \emph{selection-biased eval}, and the
bias axis is taxon abundance. After four LB attempts costing slots, our
working rule for future submissions is: target local macro deltas in
the range $[-0.05, -0.02]$ to search for LB positives; local gains
greater than $+0.05$ are reliable warning flags for rare-taxa overfit.
This is the actionable output of five local experiments (exp45a-f)
grounded in five LB submissions.

\subsection{exp47 --- Boredom-style SED faithful reproduction}
\label{sec:exp47}

\paragraph{Motivation.} The 2025 winner ``Boredom'' reported LB~0.947
with a single EfficientNet SED trained on \texttt{train\_audio} plus
labeled soundscapes only --- no Perch, no knowledge distillation, no
unlabeled pseudo-labeling. Our strongest single SED so far
(exp41f, HGNet-B0 + ensemble-teacher pseudo + labeled SS) reaches
local macro 0.902 on our rare-heavy 11-file held-out. The
reproducibility of Boredom's 0.947 number determines whether our
pipeline has an implementation gap or a fundamental ceiling.

\paragraph{Recipe.} We reproduce the consolidated 2025 top-scorer
insight:
\begin{itemize}
\item Backbone: HGNetV2-B0 (\texttt{hgnetv2\_b0.ssld\_stage2\_ft\_in1k})
\item Input: 20\,s raw audio $\to$ GPU mel-spectrogram
  ($n_{\text{fft}}=2048$, hop $512$, $n_{\text{mels}}=128$,
  $f\in[50, 14000]$\,Hz)
\item Input normalisation: \texttt{BatchNorm2d}$(n_{\text{mels}})$
  applied on the frequency dimension --- without this, fp32 logits
  overflow during warmup
\item Augmentation: raw-waveform mixup ($\alpha{=}0.5$, $p{=}0.6$) +
  SpecAugment ($F{=}16$, $T{=}40$); no background-mixing
\item Head: attention-pooled clipwise + framewise logits, hybrid loss
  BCE(clipwise) + BCE(framewise\_max), with primary-label
  one-hot and secondary-label weight 0.3
\item Optimiser: AdamW, LR $10^{-3}$, WD $10^{-2}$, cosine over 30\,ep
  with 3\,ep linear warmup, grad clip 1.0, pure fp32 (AMP overflows on
  RTX 5090 compute cap 12.0a)
\item Data: 28\,526 \texttt{train\_audio} clips (80\% stratified by
  \texttt{primary\_label}) + 617 labeled SS windows (55 files);
  weighted sampler gives $\sim$85\% train\_audio / 15\% SS per batch
\end{itemize}

\paragraph{Dual evaluation.} Two held-out AUCs per epoch, both
macro-averaged over classes with both positives and negatives:
\begin{itemize}
\item \textbf{val\_TA} (171 evaluable classes) --- 20\% of
  \texttt{train\_audio}, clean single-label, dominated by common Aves.
  Plays the role of a clean-domain LB proxy.
\item \textbf{val\_SS} (40 evaluable classes) --- the same 11 labeled-SS
  files used by exp38/41f, rare-heavy multi-label.
\end{itemize}

\paragraph{Results.} After 21 ep (run halted at patience-3 plateau on
val\_TA):
\begin{center}
\begin{tabular}{rccc}
\toprule
ep & LR & val\_TA (171 cls) & val\_SS (40 cls) \\
\midrule
01 & 0.00033 & 0.8950 & 0.7954 \\
02 & 0.00067 & 0.9342 & \textbf{0.8661} \\
05 & 0.00099 & 0.9726 & 0.8442 \\
07 & 0.00095 & 0.9764 & 0.8332 \\
11 & 0.00080 & 0.9795 & 0.8372 \\
14 & 0.00064 & 0.9852 & 0.8136 \\
16 & 0.00053 & \textbf{0.9866} & 0.7858 \\
21 & 0.00025 & 0.9866 & 0.8048 \\
\bottomrule
\end{tabular}
\end{center}

\paragraph{Finding 1: Boredom's headline is train\_audio-domain.}
A single HGNet-B0 reaches val\_TA 0.987 within 20 epochs under the
published Boredom recipe --- \emph{well above} the 0.947 claim. The
discrepancy tells us Boredom's 0.947 is almost certainly not the LB:
it is more consistent with a clean \texttt{train\_audio} held-out
macro than with the site-heterogeneous public test LB. Our own LB
ceiling of 0.929 is therefore not a ``Boredom gap'' to be closed; it
is a hidden-site covariate shift that a train\_audio-tuned model
cannot escape.

\paragraph{Finding 2: Within-run val\_TA--val\_SS anti-correlation.}
val\_TA monotonically rises $0.895 \to 0.987$, yet val\_SS peaks
early at ep02 ($0.866$) and decays to $0.786$ at ep16. This is the
selection-biased anti-correlation of Section~\ref{sec:exp45}, but
observed \emph{within a single optimisation trajectory}. The
trade-off is not a model-capacity limit. It is capacity
\emph{allocation}: BCE over 234 classes has a majority-class-dominated
gradient; polish cycles on common Aves memorise confidently wrong
responses on 27 double-blind sonotypes and other data-scarcity-cliff
taxa, driving their AUC below chance.

\paragraph{Finding 3: SED strengthening does not unlock additional
blend lever.} Perch already near-saturates common-Aves classes (LB
0.929 with no SED ensemble). Strengthening the SED component
(SED29 alone 0.737 $\to$ SED41f alone 0.878 $\to$ exp47 val\_TA 0.987,
all measured on their respective held-outs) improves
\emph{within-domain} single-model AUC on common classes, but adds no
information that Perch does not already have on those classes. The
only classes where SED could in principle contribute new
information are the rare 27 double-blind species --- and on those
classes, every SED trained on the same imbalanced data remains weak.
The SED axis is orthogonal to Perch only on rare species, and on rare
species it is orthogonally weak. This reframes the v17 LB regression
(SED29 $\to$ SED41f swap, $-0.007$ LB) as an expected outcome: the
blend was paying for marginal common-class noise without the rare-class
orthogonal signal needed to offset it.

\paragraph{Implication for the paper thesis.} The failure mode is
universal across: (i) the feature backbone (Perch, supervised
pre-trained on 14\,795 bird classes), (ii) distilled students (exp41f
family), and (iii) from-scratch SEDs (exp47, Boredom recipe). All
three cohorts achieve strong common-Aves metrics and all three
collapse on the 31 species that are absent or near-absent from
\texttt{train\_audio}. \textbf{A bigger/better SED is not the lever.}
The lever must inject rare-class information from outside the
\texttt{train\_audio} $\cup$ labeled SS distribution (external
bioacoustic datasets, synthetic physical-model audio, or
domain-adaptation signals on unlabeled test-site SS).

\subsection{exp48 --- Six-hour deep dive on error patterns and post-hoc levers}
\label{sec:exp48}

\paragraph{Motivation.} The exp47 conclusion that ``bigger SED does not
help'' rested on aggregate measures. To find remaining levers within
our data regime, we performed a structured six-hour audit on the 11-file
held-out: six new error-pattern signatures (exp48a), and five candidate
post-hoc interventions (exp48b-g).

\paragraph{Six novel error patterns (exp48a).}
\begin{itemize}
\item \textbf{Site variance is the dominant error axis.} Per-site macro
  AUC on v12: S13 0.833, S22 0.780, S23 0.668, S19 0.592, S08 0.517 ---
  a 0.316 spread.
\item \textbf{Hour-of-day matters.} 20:00-24:00 (night) macro 0.794;
  16:00-20:00 (dusk) 0.592. Diurnal species overlap dominates.
\item \textbf{Errors are confident-wrong, not low-signal.} For all
  bottom-8 classes, the top-10 negative-row prediction exceeds the
  positive-row median (e.g., 326272 Amphibia: neg-top10 0.944 vs
  pos-median 0.398). Perch does not abstain on unknown taxa --- it
  assigns them high confidence to the wrong species.
\item \textbf{Confusion is structurally clustered.} 47158son21, 22, 23
  (Insect sonotypes) all project to the same Aves triplet
  \{grhtan1, greant1, compot1\}. 25073, 67107, 326272 (Amphibia) all
  project to \{bcwfin2, rutjac1, tattin1\}. This regularity enables a
  cluster-based rewrite rule.
\item \textbf{Within-file prediction SD tracks signal quality.} For
  bottom-8 classes, within-file SD is 0.01--0.02 (flat curves --- purely
  file-level bias). For well-performing classes, SD is 0.08+ (temporal
  tracking of species calls). Flat predictions signal absence of
  acoustic information.
\item \textbf{Prediction range is compressed to [0.47, 0.51].}
  Multiplicative priors at inference have enormous ranking leverage
  because absolute confidence is near-uniform; all AUC is driven by
  small relative differences.
\end{itemize}

\paragraph{Five post-hoc levers, leak-free evaluation (exp48b--g).}
All priors and cluster definitions are derived from the 55-file train
split; evaluation on the 11-file held-out. Base: v12 macro 0.7143.

\begin{center}
\begin{tabular}{lcc}
\toprule
Lever (best config)                         & macro $\Delta$ & Aves $\Delta$ \\
\midrule
site prior, soft $\tau{=}0.75$              & $+0.117$  & $+0.006$ \\
cluster rewrite alone, $\alpha{=}2$         & $+0.034$  & $+0.003$ \\
site + cluster, $\tau{=}0.5$, $\alpha{=}2$  & $+0.118$  & $+0.004$ \\
site + cluster + exp47 blend ($w_{47}{=}0.25$, $\tau{=}0.75$, $\alpha{=}4$) & $\mathbf{+0.141}$ & $\mathbf{+0.062}$ \\
\bottomrule
\end{tabular}
\end{center}

Eval-derived cluster rewrite (leaky reference) reaches $\Delta = +0.090$;
the leak-free version ($+0.034$) is about one-third. The remaining
$+0.034$ is genuine transferable signal. Combining site prior with
cluster rewrite yields $+0.118$, showing the two levers are
non-redundant (site prior targets geographic species absence; cluster
rewrite targets within-window Aves-cluster co-firing).

\paragraph{exp48f --- modality independence audit.}
We ran exp47 inference on the 11-file held-out and compared four
teachers head-to-head.

\begin{center}
\begin{tabular}{lcc}
\toprule
Pair  & Spearman $\rho$ (preds) & residual-corr (pred-$y$) \\
\midrule
perch vs sed29          & $+0.003$ & $+0.792$ \\
perch vs sed41f         & $-0.037$ & $+0.695$ \\
\textbf{perch vs exp47} & $\mathbf{+0.011}$ & $\mathbf{+0.640}$ \\
sed29 vs sed41f         & $+0.349$ & $+0.901$ \\
sed41f vs exp47         & $+0.675$ & $+0.937$ \\
\bottomrule
\end{tabular}
\end{center}

exp47 has both the weakest signal dependence and the weakest error
dependence with Perch. The gap between exp47 and sed41f (both see
labeled SS) is attributable to exp47's absence of Perch-derived
pseudo-labels: sed41f's pseudo-distillation path imports Perch's error
structure.

\paragraph{exp48f --- bottom-8 per-teacher AUC: the mechanism.}
\begin{center}
\begin{tabular}{lccccc}
\toprule
class & taxon & Perch & sed29 & sed41f & \textbf{exp47} \\
\midrule
516975   & Mammalia & 0.500 & 0.000 & 0.000 & \textbf{1.000} \\
67107    & Amphibia & 0.357 & 0.408 & 0.490 & \textbf{0.873} \\
bafcur1  & Aves     & 0.278 & 0.593 & 0.946 & \textbf{0.966} \\
74113    & Mammalia & 0.375 & 0.498 & 0.815 & 0.615 \\
25073    & Amphibia & 0.500 & 0.414 & \textbf{1.000} & 0.434 \\
326272   & Amphibia & 0.414 & 0.438 & \textbf{0.755} & 0.442 \\
116570   & Reptilia & 0.500 & 0.441 & 0.818 & 0.646 \\
47158son11 & Insecta & 0.500 & 0.479 & \textbf{1.000} & 0.992 \\
\bottomrule
\end{tabular}
\end{center}

Two classes (516975 Mico melanurus, 67107) are solved \emph{only} by
exp47 at $0.87$-$1.00$ while Perch, sed29 and sed41f cluster around
chance. The controlled comparison --- sed41f and exp47 both see labeled
SS, same backbone family, same data budget --- isolates the
pseudo-distillation factor. 516975 is unmapped from Perch's
14,795-class taxonomy: Perch pseudo-labels on 516975's true
positive windows are confidently ``Aves-X'', which sed41f's
distillation loss treats as ground truth, overwriting the clean labeled
SS positives. exp47, trained with no pseudo supervision, learns the
clean S09 Mico melanurus labels directly.

\paragraph{Lesson refinement.} Our earlier claim ``bigger SED does not
help'' held in aggregate but hides a bimodal per-class reality:
\begin{itemize}
\item For common Aves (Perch near-saturated): more capacity adds no
  information, as asserted.
\item For rare classes in labeled SS: Perch-distillation is actively
  harmful, not neutral. Training without pseudo recovers clean-label
  signal. This is the concrete mechanism behind Boredom's (2025
  winner) ``no KD'' recipe choice.
\end{itemize}

\paragraph{Anti-correlation rule refinement.}
Across four submitted LB experiments, local-macro $\Delta$ sign
correlated inversely with LB $\Delta$ --- but the critical signature is
not the local macro itself, it is the \textbf{sign of Aves $\Delta$}.

\begin{center}
\begin{tabular}{lccc}
\toprule
Config                    & local $\Delta$ & Aves $\Delta$ & LB $\Delta$ \\
\midrule
v18 (pushes rare down Aves) & $+0.131$ & negative & $-0.013$ \\
v19 (rare overlay)          & $+0.167$ & negative & $-0.022$ \\
v20 (gate on v17)           & $-0.032$ & $\approx 0$ & $+0.007$ \\
v21 (gate on v12)           & $-0.033$ & $\approx 0$ & $\phantom{+}0.000$ \\
exp48g aggressive (proposed) & $+0.141$ & $+0.062$ & untested \\
\bottomrule
\end{tabular}
\end{center}

Interventions that raise local macro by trading Aves for rare have
historically regressed LB (v18/v19). Interventions that touch Aves
neutrally neither help nor hurt (v20/v21). exp48g is the first config
where local $+$ is achieved \emph{together with Aves $+$}, putting it
outside the previously established anti-correlation regime.

\paragraph{Damage check.} The aggressive top config has one
class with catastrophic drop: 1491113 (Amphibia) $0.565 \to 0.118$
($-0.447$) due to mis-defined cluster triggers. A per-class cluster
whitelist is recommended before LB submission.

\subsection{exp48g LB result and the v22 catastrophe}
\label{sec:v22_catastrophe}

The exp48g aggressive configuration was submitted as v22 (Kaggle kernel
v22) and \textbf{regressed by $-0.016$ LB} (0.929 v12 baseline $\to$ 0.913).
This is the largest LB regression we have observed and refutes the
``Aves $\Delta$ positive predicts LB positive'' hypothesis tabled at the
end of \S\ref{sec:exp48}.

\paragraph{Audit.} On the 66-file labeled SS held-out (eval $+$ training),
v22 macro AUC is $0.910$ vs v12 $0.722$, an apparent $+0.188$ improvement
that includes Aves at $+0.20$. Yet LB regressed $-0.016$. The
\textbf{Spearman row-correlation between v12 and v22 predictions
collapses to $0.745$} (vs $\geq 0.95$ for our other interventions),
indicating that v22 reshuffled $\sim$25\% of the within-row class
ranking on every test window. The site-prior factor $\tau \cdot
P(\textrm{species} \mid \textrm{site}) + (1-\tau)$ is the dominant
contributor: at $\tau = 0.75$ it directly multiplies the predicted
score by the per-site empirical frequency from the train soundscape
split. The cluster-rewrite $\alpha = 2$ further amplifies any class for
which the trigger-Aves triplet co-fires.

\paragraph{Mechanism.} Both site prior and cluster rewrite are
\emph{train-distribution lookup tables}. They encode regularities
specific to the 55-file train SS (sites S03-S23). On hidden-site test,
those regularities do not hold. The cluster rewrite specifically
multiplies the rare-target score by $1 + \alpha \cdot
\min(\textrm{trigger-prob})$; on the test distribution the same trigger
Aves fire without the rare target being present, generating false
positives. This shows up as: very large local AUC gain on classes
matching the cluster patterns (which dominate the local 40-class
evaluable set), large LB regression because the patterns do not
generalise.

\paragraph{Updated lever taxonomy.} After v22, our submitted-and-scored
levers split into three classes:

\begin{center}
\begin{tabular}{lcc}
\toprule
Lever class & Mechanism & LB transfer \\
\midrule
Train-SS-derived multiplicative lookup & v18, v19, v22 & $-0.013$ to $-0.022$ (always negative) \\
Self-consistency post-hoc gate         & v20, v21 V9 taxon gate    & neutral (0 to $+0.007$) \\
Base-pipeline blend modification       & v13, v14, v17 (Perch-derived SEDs) & $-0.003$ to $-0.007$ \\
\bottomrule
\end{tabular}
\end{center}

The first class is a strict negative pattern: any rule learned from
labeled-SS row-level distribution fails on hidden-site test.

\subsection{exp50 --- Boredom recipe with 2025 soundscape BG mixing}
\label{sec:exp50}

Motivated by exp47's ``no Perch pseudo'' insight (\S\ref{sec:exp47}), we
re-trained the same architecture with one critical augmentation:
during raw-waveform mixup, with probability $p_{\text{BG}} = 0.4$,
mix the source clip with a quiet 5-sec window randomly sampled from the
9~726 \emph{2025-edition} soundscapes (Colombian Upper Magdalena
Valley, geographically distinct from the 2026 Pantanal training sites).
$\alpha \sim \mathcal{U}(0.3, 0.7)$. We extract 11~920 quiet windows
from 4000 sampled 2025 soundscapes (RMS in bottom 40~\% of file,
$\textrm{abs}_{\max} < 0.95$, audible) into a 7~GB pool.

Hypothesis: the Pantanal-only 2026 SS gives the model a stable site
fingerprint to memorise; mixing with Colombian background forces the
model to discriminate species against backgrounds it has never seen at
training, encouraging a site-invariant decision boundary.

\paragraph{Result.} exp50 best epoch 16:
\begin{itemize}
\item val\_TA (170 train\_audio classes, common-Aves dominated): $0.9882$
  vs exp47's $0.9866$ ($+0.0016$, no harm to common).
\item val\_SS (40 labeled-SS held-out classes, rare-heavy): $0.8378$ vs
  exp47's $0.7858$ \textbf{($+0.0520$)}.
\end{itemize}

The 2025 BG augmentation gives a $+5\%$ rare-class gain at zero common-class
cost. Per-class on the 11-file held-out (\S\ref{sec:exp52} extension):
exp50 alone uniquely rescues 516975 (Mico melanurus, Mammalia) from $0.012 \to
1.000$, 67107 (Amphibia) from $0.357 \to 0.873$. sed41f, despite seeing
the same labeled SS, does not rescue these because Perch-pseudo
labels overwrite the labeled-SS gradient signal.

\subsection{exp50 LB integration --- breaking the v12 ceiling}
\label{sec:exp50_lb}

We tested the exp50 SED at progressively higher blend weights against
the production v26 baseline. Weight sweep:

\begin{center}
\begin{tabular}{rrr}
\toprule
$w_{\text{exp50}}$ & Config name & LB \\
\midrule
$0.10$ (with $w_{\text{SED29}} = 0.10$) & v23 & $0.928$ \\
$0.20$ (SED29 dropped)                   & v24 & $\mathbf{0.930}$ \\
$0.30$ (SED29 dropped)                   & \textbf{v26} & $\mathbf{0.931}$ \\
$0.50$ (SED29 dropped)                   & v25 & $0.928$ \\
\bottomrule
\end{tabular}
\end{center}

\textbf{v26 is the first submission to exceed the long-standing v12 LB
ceiling of 0.929.} The slope $dLB/dw_{50}$ peaks at $0.3$, regresses
beyond. Direct comparison: v17 used the same blend structure (0.8 P + 0.2
SED41f) and got LB $0.922$; substituting exp50 for SED41f lifts LB to
$0.930$ ($+0.008$). The Perch-independence (residual-correlation with
Perch on 11-file held-out: exp50 = 0.640 vs sed41f = 0.695 vs sed29 =
0.792) translates directly to LB transfer quality.

\paragraph{Class-level mechanism (exp64).} Per-class AUC change on the
66-file labeled SS, v12 $\to$ v26:

\begin{itemize}
\item 32 of 234 classes have KS distribution shift $> 0.3$. \textbf{25 of
  these are 47158son* Insecta sonotypes} ($n_{\text{train\_audio}} = 0$
  in 2026); the remaining 7 are non-Aves (Mammalia, Reptilia, some
  Amphibia). \textbf{0 of 162 Aves classes} shifted by $> 0.3$.
\item v26 lifts the bottom-AUC classes drastically: 516975 ($0.012 \to
  1.000$), plcjay1 ($0.153 \to 0.999$), 65377 ($0.416 \to 0.998$).
\item Common Aves with $v_{12} \geq 0.9$ have mean $\Delta = +0.043$
  (small uplift, no regression).
\end{itemize}

The mechanism is targeted: \textbf{exp50 fills the rare-class void
without disturbing common-class predictions where Perch already
dominates}. This is consistent with a residual-correlation hypothesis
of LB transfer: ensembles of less-correlated teachers transfer better
than ensembles of correlated teachers.

\subsection{exp64 --- Perch's structural dead zones}
\label{sec:exp64}

To understand the v26 gain mechanism precisely, we computed Perch's
$p_{99}$ (top-1\% predicted probability) across all 10~658 unlabeled
2026 soundscapes per class. This is a model property determined entirely
by Perch's training data; it does not depend on the labeled-SS
distribution.

\begin{center}
\begin{tabular}{lcc}
\toprule
Taxon & Total cls & $p_{99} \geq 0.3$ (``Perch alive'') \\
\midrule
Aves & 162 & 91 (56\%) \\
Amphibia & 35 & 11 (31\%) \\
\textbf{Insecta} & 28 & \textbf{0 (0\%)} \\
Mammalia & 8 & 2 (25\%) \\
Reptilia & 1 & 0 (0\%) \\
\midrule
Total & 234 & 104 (44\%) \\
\bottomrule
\end{tabular}
\end{center}

\textbf{Perch's $p_{99} < 0.1$ on every Insecta sonotype.} Perch's training
includes only 3 of 28 species in our Insecta set; the remaining 25
($\sim$\texttt{47158son*}) and the 11 with $p_{99} < 0.1$ in non-Aves
classes are structurally outside Perch's representational coverage.
This explains why v12 (Perch-dominated) is hard-capped on these
classes, and why exp50's signal on the same classes lifts the blend.

\subsection{exp65--66 --- Multi-criteria class-conditional routing}
\label{sec:exp65_66}

We then tested whether the Perch-aliveness signal could be turned into
a per-class blend rule. Several variants were evaluated on the 11-file
held-out:

\begin{center}
\begin{tabular}{lccc}
\toprule
Rule & local $\Delta$ & Spearman vs v26 & Aves $\Delta$ \\
\midrule
v26 (global $0.7P + 0.3 \cdot \text{exp50}$) & $0$ & $1.000$ & $0$ \\
R5 ($p_{99} < 0.1$ AND exp50\_$p_{99} \geq 0.3$ $\to$ $w_{50}=1$, $n=4$) & $+0.016$ & $0.992$ & $+0.019$ \\
R6 (continuous $w_{50} = 1 - 0.7 \cdot p_{99}$) & $+0.018$ & $0.956$ & $+0.018$ \\
R8 (3-tier: dead/semi/alive) & $+0.017$ & $0.991$ & $+0.019$ \\
R9 (drop exp50 where exp50 dead) & $\mathbf{-0.047}$ & $0.966$ & $\mathbf{-0.073}$ \\
\bottomrule
\end{tabular}
\end{center}

R9's collapse is the diagnostic surprise: \textbf{when we surgically
remove exp50 weight from the 163 classes where exp50 itself has
$p_{99} < 0.1$, local macro drops by $-0.047$.} Because exp50 produces
near-constant low-magnitude output on those 163 classes, the
$0.3 \cdot z(\text{exp50})$ term in v26 is mostly amplified random
noise --- and yet removing the noise hurts.

This recasts our understanding of what v26 actually does. The
$+0.002$ LB gain (v12 $\to$ v26) is at least partially attributable to a
\textbf{noise-regularisation effect} from the additive $0.3 \cdot
z(\text{noise})$ on Perch-strong classes, not only to the targeted
rare-class signal on Perch-dead classes. Removing the noise removes a
small but positive regulariser on the cross-class macro AUC ranking.

\paragraph{Implications.} The local $+0.016$--$+0.018$ gain of routing
schemes (R5, R6, R8) is most likely \emph{also} a partial-noise effect
on the labeled SS distribution and does not necessarily transfer to LB.
This is consistent with the v22 / v28 / v29 pattern, where each routing
attempt with apparently safe Spearman ($\geq 0.99$) and positive Aves
$\Delta$ has nonetheless regressed LB by $-0.001$ to $-0.016$.

\subsection{exp68 --- Quick rejections of Gauss / rank-norm levers}
\label{sec:exp68}

Multi-$\sigma$ Gauss ensembles ($\sigma \in \{0.3, 0.5, 0.7, 1.0\}$
averaged) and per-class quantile-to-normal rank normalisation were
tested as alternatives to the routing schemes:

\begin{itemize}
\item Multi-$\sigma$ averaging gives macro $\Delta \leq +0.0024$
  (negligible) with Spearman $\geq 0.95$.
\item Per-class quantile rank-normalisation is per-class AUC-invariant
  by construction and yields macro $\Delta = 0$.
\item Per-class std equalisation: $-0.0003$ (slight regression).
\item Per-row entropy adaptive blend: $+0.004$ (margin-of-error).
\end{itemize}

These post-hoc levers do not move the needle. Rejected as candidates.

\subsection{exp59 --- ConvNeXt SED diversity test}
\label{sec:exp59}

To probe whether architecture diversity translates to prediction
diversity, we trained a ConvNeXt-tiny SED with the same Boredom recipe
$+$ 2025 BG mixing as exp50 (HGNet-B0). Best epoch 18:
val\_TA $0.9872$ vs exp50 $0.9882$ (parity), val\_SS $0.8589$ vs
exp50 $0.8378$ ($+0.021$).

\paragraph{Diversity check on 66-file labeled SS.}
\begin{center}
\begin{tabular}{lc}
\toprule
Pair & Pearson on raw predictions \\
\midrule
Perch $\leftrightarrow$ exp50 & $-0.058$ \\
Perch $\leftrightarrow$ exp59 & $-0.058$ \\
\textbf{exp50 $\leftrightarrow$ exp59} & $\mathbf{+0.981}$ \\
\bottomrule
\end{tabular}
\end{center}

\textbf{Different architectures with the same training recipe produce
near-identical predictions.} The training procedure (data + recipe)
dominates inductive bias more than the backbone choice. ConvNeXt and
HGNet-B0 are essentially substitutable in our blend.

\paragraph{Local 4-way blend variants.} We tested various 3-teacher
($P + \text{exp50} + \text{exp59}$) blends on held-out 11:

\begin{center}
\begin{tabular}{lcc}
\toprule
Config & macro $\Delta$ vs v26 & Spearman \\
\midrule
$0.7P + 0.2 \cdot \text{exp50} + 0.1 \cdot \text{exp59}$ & $+0.005$ & $0.994$ \\
$0.7P + 0.15 + 0.15$ & $+0.009$ & $0.988$ \\
$0.6P + 0.2 + 0.2$ & $+0.019$ & $0.985$ \\
$0.6P + 0.0 + 0.4$ (drop exp50) & $+0.021$ & $0.946$ \\
\bottomrule
\end{tabular}
\end{center}

The 4-way ensemble offers marginal local gains because exp50 and exp59
are too similar. To get genuine improvement from a 4-way blend we
would need a teacher with $\rho < 0.7$ relative to exp50 (a different
training recipe, not just a different backbone).

\subsection{Updated submission timeline (2026-04-23 to 2026-04-25)}
\label{sec:submissions_v22_v29}

\begin{center}
\begin{tabular}{rlcc}
\toprule
v & Configuration & LB & $\Delta$ vs v12 (0.929) \\
\midrule
v12 & $0.8P + 0.2 \cdot \text{SED29}$ + Gauss + V9 (production)         & $0.929$ & 0 \\
v17 & $0.8P + 0.2 \cdot \text{SED41f}$                                  & $0.922$ & $-0.007$ \\
v18 & v17 + exp44c overlay                                              & $0.916$ & $-0.013$ \\
v19 & v17 + exp44g synth-aug                                            & $0.907$ & $-0.022$ \\
v20 & v17 + V9 taxon gate                                               & $0.929$ & $0$ \\
v21 & v12 + V9 taxon gate                                               & $0.929$ & $0$ \\
v22 & v12 + V9 + exp48 (site prior $\tau{=}0.75$ + cluster $\alpha{=}2$) & $0.913$ & $\mathbf{-0.016}$ \\
v23 & $0.8P + 0.1 \cdot \text{SED29} + 0.1 \cdot \text{exp50}$           & $0.928$ & $-0.001$ \\
\textbf{v24} & $0.8P + 0.2 \cdot \text{exp50}$ (drop SED29)               & $\mathbf{0.930}$ & $\mathbf{+0.001}$ \\
v25 & $0.5P + 0.5 \cdot \text{exp50}$                                   & $0.928$ & $-0.001$ \\
\textbf{v26} & $0.7P + 0.3 \cdot \text{exp50}$ (peak)                    & $\mathbf{0.931}$ & $\mathbf{+0.002}$ \\
v27 & v26 + class-conditional Aves vs non-Aves split (push)              & $0.927$ & $-0.002$ (worst case in slot 2) \\
v28 & v26 + exp51 27-head additive @ $w{=}0.10$                          & $0.929$ & $-0.002$ \\
v29 & v26 + per-class adaptive blend (R1/R3 from OOF on train SS)         & $0.927$ & $-0.004$ \\
\bottomrule
\end{tabular}
\end{center}

\paragraph{Reading the table.}
The pattern of LB regressions vs the v26 ($0.931$) baseline reveals a
strong invariant: \textbf{any modification beyond a single global blend
weight regresses LB.} Even when local Spearman vs v26 stays
$\geq 0.99$ and Aves $\Delta \geq 0$ (v28, v29), the LB outcome is
$-0.001$ to $-0.004$. The v26 weight $0.7P + 0.3 \cdot \text{exp50}$
appears to be a Pareto-frontier of the structurally simple
``Perch-primary plus Perch-independent SED'' blend.

\paragraph{What this is not.}
This is not the global LB ceiling. The leaderboard contains
submissions at $0.933$+ from competitors. Our pipeline-shape ceiling is
$0.931$; competitors at higher LB are using either (a) different probe
architectures, (b) successful per-class calibration we have not
discovered, (c) richer ensemble diversity from genuinely different
training recipes, or (d) better external data. None of these have been
ruled in or out by our experiments.

\subsection{exp67 --- EfficientNet-B0-NS SED (rejected)}
\label{sec:exp67}

A third architecture backbone (\texttt{tf\_efficientnet\_b0.ns\_jft\_in1k},
4.61M params) trained with the identical Boredom recipe + 2025 BG
mixing. Training was unstable: val\_TA fluctuated ($0.97 \to 0.83 \to
0.85$ across ep4-ep9) and early-stopped at ep9. Best epoch 4 reached
val\_TA $0.9661$ / val\_SS $0.7603$ --- significantly below
exp50 (val\_SS $0.838$) on rare classes.

Diversity check: Pearson(exp50, exp67) on 66-file labeled SS = $0.936$,
similarly correlated as exp50 $\leftrightarrow$ exp59 ($0.981$). Adding
exp67 to the v26 base (4-way blend) regresses local macro by
$-0.009$ to $-0.028$ across all weight configurations. Insecta taxon
takes the largest hit ($-0.10$ to $-0.12$) due to exp67's weak
rare-class signal.

\paragraph{Verdict.} EfficientNet-B0-NS confirmed as a weaker substitute
than HGNet-B0 in our recipe (consistent with Salman's 2025-discussion
note that EffNet-NS underperforms HGNet for this task). Rejected as a
4-way ensemble member.

\subsection{Five-option lever audit summary (2026-04-25)}
\label{sec:five_lever_audit}

After v26 ($0.931$ LB), we systematically validated five candidate
levers locally before considering a submission slot.

\begin{center}
\begin{tabular}{lll}
\toprule
Option & Local (held-out 11) effect & Verdict \\
\midrule
1. R5 per-class routing (aliveness-based) & $+0.016$ macro, sp $0.992$ & possibly LB-neutral by noise hypothesis \\
2. Probe arch swap (PCA / LogReg C / MLP) & not tested locally          & deferred (notebook OOF retrain heavy) \\
3. Multi-$\sigma$ Gauss ensemble          & $+0.0024$ best (negligible)& reject (no movement) \\
4. Test-time rank normalisation           & $0$ (AUC-invariant)        & reject (math) \\
5. EffNet-B0-NS additional teacher        & $-0.009$ to $-0.028$       & reject (weaker than exp50, Pearson 0.94 with exp50) \\
\bottomrule
\end{tabular}
\end{center}

\paragraph{Conclusions.}
\begin{itemize}
\item Three of five options (3, 4, 5) demonstrably do not help locally.
  None are worth a Kaggle submission slot.
\item Option 1 (R5 routing) shows local $+0.016$ but the
  noise-regularisation hypothesis predicts LB neutral-or-negative
  transfer (consistent with v22, v28, v29 each showing local positive
  with LB regression).
\item Option 2 (probe arch) remains an open lever but requires
  significant notebook refactoring; deferred until a clean window
  exists.
\end{itemize}

\paragraph{Practical implication.} The local audit returned no
structurally robust LB-gain candidate beyond v26. We retain v26
($0.931$) as the production baseline and defer further submissions
until a genuinely new lever is identified --- e.g., a different
training recipe (not just architecture) producing a low-correlation
($\rho < 0.7$ with exp50) teacher, or external bird/non-bird audio
data extending coverage of Perch-dead classes.

\subsection{v32 R5 aliveness routing --- final test of class-conditional levers}
\label{sec:v32}

After the local audit (\S\ref{sec:exp65_66}) showed R5 (Perch-dead AND
exp50-alive, $n=4$ classes) was the only candidate combining (a)
sp\_row $\geq 0.99$ vs v26, (b) local macro $\Delta > 0.01$, and (c)
\emph{universal model property} as the routing signal (Perch's $p_{99}$
across 10K unlabeled SS, which depends on Perch's training data, not on
the labeled-SS distribution), we submitted v32 = v26 + R5 routing on the
4 classes (22961, 22967, 326272 Amphibia + whtdov Aves).

\textbf{LB result: 0.930 ($-0.001$ vs v26).}

\paragraph{Updated lever taxonomy after v32.}

\begin{center}
\begin{tabular}{lcc}
\toprule
Routing/lookup signal source             & Submission   & LB $\Delta$ \\
\midrule
Train-SS row patterns (cluster co-firing) & v22          & $-0.016$ \\
Per-class OOF AUC (train-distribution)    & v29          & $-0.004$ \\
Direct dedicated-model output             & v28          & $-0.002$ \\
\textbf{Model property ($p_{99}$ on 10K unlabeled)} & \textbf{v32} & $\mathbf{-0.001}$ \\
\bottomrule
\end{tabular}
\end{center}

\textbf{All four mechanisms regress LB.} The structural-distance ranking
predicted v32 should be the safest: $p_{99}$ is computed from
unsupervised Perch outputs on test-domain audio and reflects Perch's
training coverage, which is determined entirely by what Perch was
trained on (not by our labeled SS or any data-derived signal). v32 is
the most surgical (only 4 columns altered) and most universal-property
based intervention we could design. It still regressed.

\paragraph{Conclusion: v26 is the local Pareto-frontier.}
The v26 configuration $0.7 \cdot \mathrm{Perch} + 0.3 \cdot
\mathrm{exp50}$ + Gauss + V9 gate is the global optimum of our pipeline.
Any per-class modification --- multiplicative or additive, supervised or
universal-property based --- regresses LB. The likely mechanism is
\textbf{noise-regularisation}: the $0.3 \cdot z(\mathrm{exp50}_\mathrm{dead})$
component on Perch-strong classes acts as a small ranking regulariser
that any precise routing attempts to ``correct'' away. R9 from
exp66 ($-0.047$ local when removing exp50 weight on exp50-dead
classes) supports this --- the noise contribution is positive on local
macro, and presumably also on LB.

\paragraph{Path beyond v26 (untested).}
Given our pipeline-shape ceiling at 0.931 confirmed via four LB
attempts (v22, v28, v29, v32), further LB gains require structural
departures from the current pipeline:
\begin{itemize}
\item Replacing or augmenting Perch with a foundation model whose
  training coverage differs (e.g., AudioMAE, BEATs, or an ensemble of
  bird foundation models).
\item Training a recipe-divergent SED (focal loss, contrastive
  objective, different audio representation) that achieves $\rho < 0.7$
  with exp50 on raw predictions, providing genuine ensemble
  diversity rather than 0.94+ correlation seen in HGNet/ConvNeXt/EffNet.
\item Stacking with strong regularisation (LightGBM or per-class linear
  with cross-validated weights), accepting the risk that train-SS
  distribution dependence may introduce v22/v29-style transfer failure.
\item External data (xeno-canto, iNaturalist) feeding a new SED with
  per-species coverage Perch lacks.
\end{itemize}

None of these have been ruled in or out. The current state of evidence
suggests v26 (0.931 LB) is our pipeline optimum, and further
experimentation should target the four structural directions above
rather than additional post-hoc levers.

\subsection{exp70 --- Row-level failure mode audit and masking tests}
\label{sec:exp70}

A direct-evidence audit of why v26 fails on its weakest classes. For
each class with v26 held-out AUC $< 0.7$, we examined positive-row
prediction magnitudes from Perch and exp50, and tested four
specific masking rules.

\paragraph{Failure mode taxonomy.} Of 12 weakest v26 classes, 7 have
\textbf{both Perch and exp50 with self-prediction $\approx 0$ on
positive rows} (\emph{both weak}). Examples:

\begin{center}
\begin{tabular}{lcccc}
\toprule
class & taxon & Perch self-mean & exp50 self-mean & v26 AUC \\
\midrule
47158son06 & Insecta & 0.50 & \textbf{0.00} & 0.07 \\
47158son17 & Insecta & 0.50 & \textbf{0.00} & 0.27 \\
326272     & Amphibia & 0.00 & 0.00 & 0.40 \\
67107      & Amphibia & 0.48 & 0.00 & 0.40 \\
74113      & Mammalia & 0.11 & 0.00 & 0.42 \\
bafcur1    & Aves     & 0.30 & 0.00 & 0.54 \\
\bottomrule
\end{tabular}
\end{center}

For the worst 7 classes, the magnitude of exp50 prediction on the
positive rows is essentially zero. This contradicts the apparent
val\_SS macro of 0.84 because per-class ranking can still be partially
informative even when magnitudes collapse. \textbf{Critically: no
masking or routing can fix a class where neither model has acoustic
signal.}

\paragraph{Masking tests.}
\begin{center}
\begin{tabular}{lcc}
\toprule
Rule (applied per cell or per row)                        & cells affected & macro $\Delta$ \\
\midrule
MASK1: Perch $> 0.5$ AND exp50 $< 0.1$ $\to$ use exp50    & 64\,828 (37\%) & $-0.027$ \\
MASK2: row-level top-Aves suppress where exp50 top is non-Aves & 615 & $-0.003$ \\
\textbf{MASK3: zero Perch on 19 $p_{99}{<}0.1$ classes}            & 19 cls         & $\mathbf{+0.016}$ \\
MASK4: top-3 Aves suppress where exp50 disagrees ($<0.1$)  & 2\,207 cells & $-0.004$ \\
\bottomrule
\end{tabular}
\end{center}

MASK3 is functionally equivalent to v32's R5 routing (LB-tested
$-0.001$). MASK1 (broad threshold-based) catastrophically breaks the
noise-regularisation balance. MASK2/MASK4 narrow rules give marginal
loss. \textbf{None of the masking schemes generates a new gain channel
beyond MASK3 = v32 = R5 already submitted to LB.}

\paragraph{Implication for path forward.}
The 12 worst classes' failure mechanism is data-level rather than
post-hoc-correctible. To improve them we need either:
\begin{itemize}
\item External labelled audio (xeno-canto for the 7 classes with
  $n_{\text{train\_audio}} < 50$), feeding a new SED.
\item TTA: multiple temporal crops averaged, recovering the marginal
  signal that exp50's single-crop misses.
\item Stacking with rich sequential / cross-class features that allow
  the meta-learner to exploit weak per-class signal.
\item Recipe-divergent SED training (focal loss, contrastive) that
  forces the model to extract rare-class signal even with sparse
  positives.
\end{itemize}

Pure post-hoc correction on the existing v26 outputs has reached its
limit; further LB gains require structural changes to the input
modality, training data, or training objective.

\subsection{exp71-73 --- External data download and fine-tune (negative)}
\label{sec:exp71_73}

To test whether external labelled audio could fix the 12 v26 worst classes
(\S\ref{sec:exp70}), we attempted limited downloads from xeno-canto and
iNaturalist for 5 priority species (Insecta sonotypes excluded due to
generic \texttt{inat\_taxon\_id} 47158 covering all of order Insecta,
no usable specific source).

\paragraph{Download outcome.}
\begin{itemize}
\item xeno-canto v2 API deprecated; v3 requires API key signup. Skipped.
\item iNaturalist research-grade observations: 43 audio recordings
  across 5 species (326272, 67107, 74113, bafcur1, litnig1).
\item After ffmpeg conversion to ogg 32kHz mono: 30 usable clips.
\item Total disk: 33 MB.
\end{itemize}

\paragraph{Quality diagnostic on raw external clips.} Running exp50 on
these 30 clips: per-clip self-prediction is bimodal --- some clips
hit $0.9+$ (clean call captured) but many sit below $0.05$ (noisy or
non-target audio dominating). Filtering by self $> 0.3$ keeps ~25
clips.

\paragraph{Fine-tune setup.} Warm-start from exp50 best ckpt; train 8
epochs at LR $5\!\times\!10^{-5}$ on combined data: 2026 train\_audio
(28.5K) + labeled SS train (617) + 25 filtered external clips
($5\times$ oversample). Held-out 11-file val\_SS evaluation each epoch.

\paragraph{Result table (ep4 best vs baseline).}
\begin{center}
\begin{tabular}{lccccc}
\toprule
species & baseline (exp50) & after ft (ep4) & $\Delta$ \\
\midrule
326272   & 0.536 & 0.586 & +0.050 \\
67107    & 0.790 & 0.740 & $-0.050$ \\
74113    & 0.765 & 0.677 & $-0.088$ ($n_{\text{pos}}=2$) \\
bafcur1  & 0.818 & 0.604 & $-0.214$ \\
litnig1  & 0.994 & 0.989 & $-0.005$ \\
\midrule
val\_SS macro (40 cls) & 0.838 & 0.851 & $+0.013$ \\
\bottomrule
\end{tabular}
\end{center}

The val\_SS macro increased by $+0.013$, but \textbf{4 of 5 target
species regressed}. The macro gain comes from non-target classes ---
i.e., a regularisation-like artefact, not a real lift on the species we
intended to help.

\paragraph{Interpretation.} 25 iNaturalist clips per species are too
few and too acoustically diverse (smartphone recordings vs labeled
SS handheld field captures) to provide a coherent positive signal for
the target species. The 2025 winners' ``data refinement'' work likely
involves manual curation, not unfiltered API downloads. Without
substantial human-in-the-loop label cleaning, a small external data
fine-tune offers no reliable gain.

\paragraph{Verdict.} External data via API alone is not the lever.
Future external-data work needs at minimum (a) per-clip manual
curation, (b) larger volume per species (200+), (c) acoustic
domain-matching to soundscape recordings. None feasible in our
remaining time budget.

\subsection{v33 --- L2c file-coherence breakthrough}
\label{sec:v33}

After 5 consecutive failures of post-hoc levers above v26 (v22, v28, v29,
v32, plus exp73 fine-tune), we tested a structurally different
intervention: \textbf{physics-based file-level temporal coherence}.

\paragraph{Lever specification.} For each test file with $N_W=12$
contiguous 5-second windows producing prediction matrix $P \in
\mathbb{R}^{12 \times 234}$, define $P_{\max}^{f,c} = \max_t P_{t,c}$.
Replace each row $t$ with
$P_{t,c}^{\text{new}} = (1-\alpha) P_{t,c} + \alpha P_{\max}^{f,c}$,
$\alpha = 0.10$.

\paragraph{Justification.} If species $c$ is genuinely present in
file $f$, it should be detected in some windows. Then
$P_{\max}^{f,c}$ reflects the highest detection confidence in $f$.
Boosting all 12 windows by $\alpha P_{\max}^{f,c}$ encodes the
\emph{physical} prior that bird/insect calls are temporally clustered:
detection in one window is evidence of presence throughout adjacent
context.

This is wider context than $\sigma{=}0.5$ Gaussian smoothing
(window 5 only sees windows 3-7 with weight $0.6$ at distance 2). The
Gaussian smooths locally but keeps file-wide max signal compressed.

\paragraph{Critical: this is NOT a train-distribution lookup.}
The transformation depends only on test predictions in the same file.
No labeled SS data, no per-class rule, no row-level pattern. Universal
physics: \emph{species presence is temporally extended}.

\paragraph{Local validation.}
\begin{center}
\begin{tabular}{lccc}
\toprule
$\alpha$ & macro $\Delta$ vs v26 & Spearman row & Aves $\Delta$ \\
\midrule
$0.10$ & $+0.0034$ & $0.994$ & $+0.007$ \\
$0.20$ & $+0.0056$ & $0.980$ & $+0.014$ \\
$0.30$ & $+0.0070$ & $0.956$ & $+0.022$ \\
\bottomrule
\end{tabular}
\end{center}

Picked $\alpha = 0.10$ for sp\_row $\geq 0.99$ guarantee.

\paragraph{LB result: 0.932.}
\textbf{First positive LB transfer above v26}. $+0.001$ vs v26 (0.931),
$+0.003$ vs v12 (0.929). The entire 5-experiment run of post-hoc
multiplicative or routing-based levers regressed; the first
physics-based intervention achieved positive transfer.

\paragraph{Updated lever taxonomy after v33.}

\begin{center}
\begin{tabular}{lcc}
\toprule
Lever class                                    & Examples         & LB transfer \\
\midrule
Train-SS row-pattern lookup                    & v22 (site+cluster) & $-0.016$ \\
OOF AUC routing                                & v29              & $-0.004$ \\
Direct dedicated-model output                  & v28              & $-0.002$ \\
Universal model property                       & v32 (Perch $p_{99}$) & $-0.001$ \\
\textbf{Physics-based (file-level coherence)}  & \textbf{v33}     & $\mathbf{+0.001}$ \\
\bottomrule
\end{tabular}
\end{center}

\paragraph{Generalisation hypothesis.}
\textbf{Universal physical priors transfer; train-distribution patterns
do not.} The hierarchy from purely sample-specific (v22) to purely
universal-physics (v33) tracks LB transfer monotonically. This points to
a class of further interventions:
\begin{itemize}
\item Cross-window temporal Markov / HMM (state-space at logit level
  with prior learned from training, but transition probabilities are
  universal physical priors).
\item Frequency-domain physical priors (insect stridulation has narrow
  bandwidth, bird calls have characteristic harmonics).
\item Ecological co-occurrence priors derived from biome physics, not
  data lookup.
\end{itemize}

These are the directions that should produce additional positive LB
gains, distinct from the failed routing/lookup family.


\section{v34 mel-iVAE z-kNN: structural failure of train-SS-fitted levers (2026-04-25)}

\subsection{Motivation}
After v33's positive LB transfer (file-max coherence, $+0.001$), we
investigated whether \emph{independent feature representations} could provide
a complementary signal targeting the 27 unmapped species (25 Insecta sonotypes
+ 2 Amphibia not in Perch's 14k-class vocabulary). Perch \emph{embeddings} on
those rows are still informative; Perch's species \emph{output heads} for
unmapped species are constant. An iVAE over raw mel might bypass this output
bottleneck.

\subsection{Setup}
Pooled mel-spectrogram (T=16 temporal × F=128 mel) on 739 labeled SS rows
(617 train / 122 eval, file-stratified). Train an auxiliary-conditioned VAE
with site one-hot + hour as auxiliary; encoder produces a 32-d $z$.
Per-class centroid in $z$-space from train positives ($\geq 3$ examples
yields 55/234 valid centroids). At inference time, cosine similarity to each
valid centroid produces a $z$-kNN score; sigmoid sharpening with factor 5
maps to $[0,1]$. Final blend with $w_z=0.05$ on the v33 output.

\subsection{Local validation looked positive}
On 122 held-out eval rows:
\begin{itemize}
\item macro AUC $\Delta = +0.026$ vs v33,
\item per-row Spearman vs v33 = $0.990$ (extremely high rank preservation),
\item Aves $\Delta +0.008$ (anti-correlation rule v2 predicts neutral/positive transfer),
\item Insecta $\Delta +0.068$, Reptilia $\Delta +0.104$.
\end{itemize}

\subsection{LB result}
\textbf{v34 LB = 0.916, $\Delta = -0.016$ vs v33 (0.932).} Largest single-step
regression in the v22--v34 series.

\subsection{Mechanistic post-mortem (exp79--80)}
A six-experiment teardown of the iVAE artifact (exp79 pilot, exp79b eval-only
re-train, exp79c unlabeled probe, exp80a site-holdout, exp80b big-pool
training, exp80c taxon classifier comparison) isolated the failure mechanism.

\paragraph{Finding 1: Site holdout collapses Insecta detection.}
Train iVAE excluding all rows from S19, then evaluate Insecta detection on
S19's held-out windows: AUC = $0.073$ (sub-random). The trained model places
S19 windows \emph{further} from Insecta centroids than non-Insecta windows,
because the encoder never saw S19's acoustic environment during training.

\paragraph{Finding 2: Same-site eval AUC is site-fingerprint, not species
acoustics.}
\begin{center}
\begin{tabular}{lcc}
\toprule
Taxon & Same-site eval AUC & LOSO site CV AUC \\
\midrule
Aves      & 0.860 & 0.611 \\
Amphibia  & 0.934 & 0.554 \\
Insecta   & \textbf{0.973} & \textbf{0.295} \\
Mammalia  & 0.614 & 0.624 \\
Reptilia  & 0.778 & 0.566 \\
\bottomrule
\end{tabular}
\end{center}

The Insecta column shows a 0.68-AUC gap between same-site and
leave-one-site-out evaluation. Same-site eval reports near-perfect detection;
LOSO is sub-random ($<0.5$) -- the classifier is \emph{worse} than chance on
held-out sites.

\paragraph{Finding 3: iVAE does not add to Perch.}
Perch+iVAE-z concat MLP achieves $0.963$ Insecta same-site AUC, slightly
\emph{below} Perch alone ($0.973$). LOSO is also slightly worse
($0.205$ vs $0.295$). Bigger iVAE training pool (30k unlabeled rows) helps
Aves detection (+$0.10$) but hurts Amphibia and Insecta -- a wider pool
dilutes the centroid signal that was already site-fingerprint, not species
acoustics.

\paragraph{Finding 4: Perch's species-output blindness is not the
informative axis.}
Perch's 14k-class species head outputs constant 0 for the 25 Insecta
sonotypes. Yet a fresh MLP on Perch's 1536-d \emph{embedding} achieves
$0.973$ Insecta same-site AUC -- the embedding contains
Insecta-discriminative structure. The bottleneck is not that Perch ``can't
hear insects''; it is that all our labeled Insecta examples come from 4
sites, so any classifier fits the site, not the insects.

\subsection{Generalised lesson}
Refines the post-v33 lever taxonomy with a fifth, deeper failure mode:

\begin{center}
\begin{tabular}{lcc}
\toprule
Lever class                                & Example         & LB \\
\midrule
Train-SS row-pattern lookup                & v22             & $-0.016$ \\
OOF AUC routing                            & v29             & $-0.004$ \\
Direct dedicated-model output              & v28             & $-0.002$ \\
\textbf{Train-SS-fitted feature-extractor + cosine kNN} & \textbf{v34} & $\mathbf{-0.016}$ \\
Universal model property                   & v32             & $-0.001$ \\
Physics-based (file-level coherence)       & v33             & $+0.001$ \\
\bottomrule
\end{tabular}
\end{center}

\textbf{Any lever derived from a structure fitted to our 55-file labeled SS
train split is poisoned by site shortcut, regardless of architecture, weight,
or rank-preservation safety.} v34 confirmed this explicitly: $w_z = 0.05$
(tiny weight), Spearman $0.990$ (preservation of $99\%$ of v33's ranking), and
local validation predicting positive transfer all failed.

The only train-SS-touching mechanism that has ever transferred positively is
exp50, which used SS data only as one input stream of a multi-task SED with
\emph{2025 soundscape background mixing} for active site-invariance training.
Pseudo-labeling, post-hoc lookup, taxon gates, and feature-knn approaches
all failed.

\subsection{Implications for the working-note paper}
We treat exp79--80 as a complete \emph{negative-result thesis} on identifiable
latent representation models for foundation-model-pretrained bioacoustic
features. The same hierarchy emerges as in our earlier exp43 chain on Perch
embeddings: pre-trained models supervised on diverse audio already encode the
species-discriminative structure that an iVAE/iVDFM would aim to recover; the
remaining variance in $z$-space is dominated by site/environment factors,
which are precisely the directions an unsupervised disentanglement objective
would either keep (insufficient identifiability) or transfer to the
auxiliary's predicted prior (insufficient disentanglement). Either way, when
the labeled subset has site-conflated taxon distributions, downstream species
or taxon detection inherits the site shortcut.

The negative result has independent value: it characterises a regime
(supervised foundation embedding + site-conflated downstream labels) in which
identifiability machinery cannot recover species structure, and points the
remaining levers toward architectures that \emph{actively suppress} site
fingerprint at training time (mixup with geographically distinct
backgrounds, domain-adversarial losses) rather than disentangle it post hoc.


\section{Q1--Q4 systematic local audit (exp82--85, 2026-04-26)}

After the v34 mel-iVAE failure refuted further iVAE-fitted lever exploration,
we audited four hypothesised lever families that we had partially explored
but never rigorously measured against the v33 LB-positive baseline. Each
hypothesis was tested with a controlled local pipeline matching the v33
production recipe (Perch + SED blend $\to$ V9 taxon gate $\to$ file-max
coherence $\alpha=0.10$), evaluated on 122 held-out file-stratified eval
rows.

\subsection{Q1: Architectural diversity (ConvNeXt-tiny)}
exp59 ConvNeXt-tiny SED was trained with the same 2025-BG recipe as exp50.
Despite distinct backbone, predictions correlate at Pearson 0.981 with
exp50.

Yet swapping into 4-way blend (Perch + exp50 + exp59) at $0.6/0.2/0.2$
weight delivers macro $\Delta = +0.0093$, Aves $\Delta = +0.017$,
sp\_row $0.999$ vs v33. A control test (single-SED at the same total weight,
$0.6P + 0.4\,\mathrm{exp50}$) yields only $+0.0035$ macro, isolating the
architecture-diversity contribution at $+0.006$ macro and $+0.008$ Aves
beyond pure weight reallocation.

Despite the high correlation, the residual 2\% information is genuinely
complementary. Architecture diversity contributes a real but moderate
boost.

\subsection{Q2: Recipe diversity (focal loss)}
exp83 substituted focal loss ($\gamma=2$, $\alpha=0.25$) for BCE on the
exp50 recipe (HGNet-B0 + 2025 BG, all other settings identical). Trained
20 epochs on RTX 5090.

Best epoch (ep 18): val\_TA $0.985$, val\_SS $0.827$. Reference exp50:
val\_TA $0.988$, val\_SS $0.838$. Focal loss alone is slightly worse
than BCE on both metrics ($-0.011$ on val\_SS).

In 4-way blend ($0.6P + 0.15\,\mathrm{exp50} + 0.15\,\mathrm{exp83}$),
macro $\Delta = +0.007$, Aves $\Delta = +0.004$, but Reptilia
$\Delta = -0.047$. In 5-way blend including focal: Reptilia drops $-0.036$.

Recipe diversity via loss-function swap does not help. Focal loss with
$\gamma = 2$ on the same data + architecture produces predictions that are
almost identical (Pearson $0.981$ with exp50) and slightly worse on rare
taxa.

\subsection{Q3a: Small external pool (exp73 with 21 clips)}
exp73 was a fine-tune of exp50 on 21 quality-filtered iNaturalist clips
across 5 target species. Substituting exp73 into the 4-way blend yields
$+0.010$ macro, Aves $+0.014$, sp\_row $0.999$ -- a marginal improvement
over Q1 (+0.001 macro). Useful but small.

\subsection{Q4: Expanded external pool (exp84b with 545 clips)}
Downloaded 788 new iNaturalist clips across 56 target species (Amphibia
35, Mammalia 8, Insecta 3 named species, Reptilia 1, plus 11 weak Aves)
via exp84\_q4\_external\_expand. Combined with existing 43 clips and
filtered by exp50 self-prediction $> 0.3$, retained 545 clips.

Fine-tuned exp50 for 12 epochs (LR $10^{-4}$, larger pool justifies
slightly higher LR). Best ep 7: val\_SS $0.8628$ vs baseline $0.8378$,
$\Delta = +0.025$ -- nearly $2\times$ the exp73 gain ($+0.013$).

Per-target progress vs exp73 baseline:
\begin{itemize}
  \item bafcur1: $0.60 \to 0.79$ (+0.19 from external clips)
  \item 67107: $0.74 \to 0.87$ (+0.13)
  \item 74113: $0.68 \to 0.77$ (+0.09)
  \item 326272: $0.59 \to 0.48$ (\textbf{regressed} despite added clips)
  \item litnig1: $0.99 \to 1.00$ (saturated)
\end{itemize}

Substituting exp84b into the 4-way blend (Q4-blend) yields macro $+0.006$,
Aves $+0.011$, sp\_row $0.999$. Combining exp84b with exp50 + exp59 (5-way)
yields macro $+0.013$, Aves $+0.024$, all five taxa positive
(Reptilia $+0.028$, Amphibia $+0.014$).

\subsection{Aggregate ranking (122 eval rows, anti-correlation rule v3)}
\begin{center}
\begin{tabular}{lcccc}
\toprule
Config                                          & macro $\Delta$ & sp\_row & Aves $\Delta$ & Predicted LB \\
\midrule
6-way: 0.5P + all 5 SEDs each 0.10              & $+0.015$       & 0.998   & $+0.024$      & A (some Reptilia drop) \\
\textbf{5-way: 0.5P + e50/e59/e73/e84b each 0.125} & $\mathbf{+0.013}$ & $\mathbf{0.999}$ & $\mathbf{+0.024}$ & \textbf{A (all taxa $\geq 0$)} \\
6-way conservative: 0.6P + each SED 0.08        & $+0.012$       & 0.999   & $+0.012$      & A \\
Q3a best: 0.6P + 0.15 e50 + 0.15 e59 + 0.10 e73 & $+0.010$       & 0.999   & $+0.014$      & A \\
Q1 best: 0.6P + 0.2 e50 + 0.2 e59               & $+0.009$       & 0.999   & $+0.017$      & A \\
Q4-blend: 0.6P + 0.2 e50 + 0.2 e84b             & $+0.006$       & 0.999   & $+0.011$      & A \\
v33 ref (0.7P + 0.3 exp50)                       & $0$            & $1.000$ & $0$           & reference \\
Q2 swap: 0.7P + 0.3 exp83 (focal alone)         & $+0.003$       & 0.998   & $-0.002$      & B (focal hurts Aves) \\
\bottomrule
\end{tabular}
\end{center}

\paragraph{Findings.}
\begin{itemize}
  \item Q1 (architecture diversity): real but modest, $+0.006$ macro beyond
    weight reallocation. Predictions are 0.98 correlated with exp50; the
    residual 2\% is what matters.
  \item Q2 (recipe diversity, focal loss): \textbf{negative}. Focal loss
    alone is slightly worse than BCE; in blend it hurts Reptilia.
  \item Q3a (small external pool): marginal $+0.001$ over Q1 alone.
  \item Q4 (large external pool, 545 filtered clips): the strongest single
    contributor. Combining $e50 + e59 + e84b$ at $0.125$ each gives all
    five taxa positive Aves, the only such configuration.
\end{itemize}

The 5-way blend is the highest-confidence LB candidate after v33. Predicted
LB transfer per anti-correlation rule v3 (Aves $+0.024$, sp\_row $0.999$,
all taxa $\geq 0$): $+0.002$ to $+0.004$ over v33 (LB 0.932), targeting
$\sim 0.934$--$0.936$.

\paragraph{What this rules out.}
Recipe diversity via loss substitution alone does not produce useful
ensemble component. Pearson correlation of $0.97$--$0.99$ across all five
SED variants (independent of recipe / architecture / external data) shows
that downstream predictions converge to a narrow function class on this
pool, regardless of training-time tweaks. Future ensemble strengthening
likely requires \emph{architecturally distinct teacher families} (e.g.,
non-CNN, Perch v3, transformer-based bioacoustic models) rather than
loss/data tweaks on the same recipe.


\section{v36 5-way blend LB failure (2026-04-26)}

The 5-way blend recommended by Q1--Q4 audit (\S exp82--85) was submitted
as v36: $0.5\,P + 0.125\,(\mathrm{exp50} + \mathrm{exp59} +
\mathrm{exp73} + \mathrm{exp84b})$, with V9 taxon gate and file-max
coherence ($\alpha = 0.10$). Local prediction by anti-correlation rule v3
classified this as ``A (likely positive)'': sp\_row $0.999$, Aves
$\Delta = +0.024$, all-five-taxon $\Delta \geq 0$, macro $\Delta = +0.013$
on 122 held-out eval rows.

\textbf{LB result: $0.915$, $\Delta = -0.017$ vs v33 $(0.932)$.}

\subsection{Anti-correlation rule v3 invalidation}
This is the second consecutive submission predicted ``class A'' by
anti-correlation rule v3 that regressed by approximately $-0.017$ LB
(v34 mel-iVAE: $-0.016$; v36 5-way: $-0.017$). The rule, which previously
worked for v33 (small Aves~$\Delta$, sp\_row $\geq 0.99$, universal
mechanism), now appears systematically wrong on a different class of
modification.

\subsection{Mechanism}
v36 differs from v33 in two coupled ways: (a) Perch weight reduced from
$0.7$ to $0.5$; (b) SED side weight increased from $0.3$ to $0.5$ across
four highly correlated SED variants (Pearson $0.97$--$0.99$ across all
pairs). The Q1--Q4 hypothesis was that the residual $1$--$3\%$
non-correlation between SEDs supplies complementary information justifying
the redistribution. The v36 LB result rules out this hypothesis: the
``ensemble diversity'' is not enough to compensate for the loss of Perch's
multi-region backbone signal at $0.5$ weight.

The Aves $\Delta = +0.024$ local was a same-site evaluation artifact:
amplifying site-conflated SED signals on labeled SS rows looks like
diversity benefit because the site fingerprints align between train and
eval splits. On hidden-site test rows, the same amplification surfaces as
worse rankings because Perch's site-invariant predictions are getting
suppressed.

\subsection{Refined invariant}
The LB-protective invariant after v22, v28, v34, v36 is sharper than
``Aves $\Delta \geq 0$ + sp\_row $\geq 0.99$'': it appears to be
\emph{$W_{\mathrm{Perch}} \geq 0.7$ in any final blend}. Every submission
that reduced $W_{\mathrm{Perch}}$ below $0.7$ has regressed; every
positive transfer (v23 0.928, v24 0.930, v26 0.931, v33 0.932) preserved
$W_{\mathrm{Perch}} \in [0.7, 0.8]$. This is consistent with our
2025-BG-trained SED (exp50) being the only train-SS-touching component
that produces site-invariant predictions; even matching its recipe
(exp59/73/83/84b) does not yield models with comparable LB-transfer
quality, because the residual differences between them re-express the
underlying site-fingerprinted feature space.

\subsection{Lessons for remaining lever space}
\begin{itemize}
  \item Future ensemble proposals must keep $W_{\mathrm{Perch}} \geq 0.7$
    in the renormalized blend, or include a non-train-SS-fitted teacher
    (e.g., a future Perch v3 or transformer-based foundation model) to
    occupy the $0.7$+ weight slot.
  \item Local same-site eval is now a confirmed-misleading metric for
    LB transfer of this lever family. We need an additional eval signal
    that explicitly tests Perch-share preservation, not just rank
    preservation.
  \item Q1--Q4 audit was a useful negative result: it definitively rules
    out incremental improvements via SED-recipe / SED-data tweaks within
    the labeled-SS pool. The next lever space is genuinely
    architecturally distinct teachers, not pipeline tweaks.
\end{itemize}


\section{Honest decomposition: how the 0.929 → 0.932 jump actually happened}
\label{sec:lb-jump-decomposition}

After v34 and v36 LB regressions reset our understanding of what
transfers, this section re-analyses the only LB-positive sequence we have
produced (v12 → v24 → v26 → v33) to identify the actual invariants. The
claim ``v33 LB 0.932 is the result of three independent +0.001 steps,
each adding a new source of site-invariance to the prediction pipeline''
is supported below.

\subsection{The three +0.001 transitions}

\begin{center}
\begin{tabular}{lllr}
\toprule
Transition & Config change                        & Mechanism                                & LB \\
\midrule
v12 (ref)  & $0.8P + 0.2$ SED29 + V9 + Gauss      & --                                       & 0.929 \\
v12 → v24  & SED29 → exp50 (weight kept at 0.2)   & training-time site invariance via 2025 BG  & 0.930 (+0.001) \\
v24 → v26  & weight 0.2 → 0.3                     & optimal mixing of two site-invariance sources & 0.931 (+0.001) \\
v26 → v33  & + file-max coherence ($\alpha = 0.10$) & universal physics: temporal extension     & 0.932 (+0.001) \\
\bottomrule
\end{tabular}
\end{center}

The total +0.003 LB was obtained from three independent steps of equal
size, each contributing a distinct new mechanism. There was no single
dominant lever.

\subsection{Step 1 mechanism (v12 → v24): training-time site invariance via 2025 BG mixing}

SED29 was a HGNet-B0 trained on \emph{train\_audio only}, no labeled-SS,
no BG mixing. exp50 is the same architecture with the Boredom recipe
(20-second clips, raw waveform mixup, BCE on clip+framewise max) plus a
critical addition: with $p = 0.4$, the mixup partner is a quiet 5-second
window from the 2025 soundscape pool (Colombian Upper Magdalena, a
geographically distinct region from the Pantanal source of our 2026 data),
blended at $\alpha \sim U(0.3, 0.7)$.

This forces the SED to produce predictions invariant to the acoustic
environment of the 2026 site fingerprints, because training labels stay
fixed when the BG region changes. The downstream effect on prediction
diversity is measurable. On the 11-file held-out, residual-correlation of
each SED with Perch (after centring on per-class label):

\begin{center}
\begin{tabular}{lc}
\toprule
SED variant & residual-corr(Perch, SED) \\
\midrule
SED29 (train\_audio only)               & 0.792 \\
SED41f (Perch-pseudo distilled)         & 0.695 \\
exp47 (Boredom recipe, no 2025 BG)      & 0.640 \\
\textbf{exp50 (Boredom + 2025 BG)}       & \textbf{0.640} \\
\bottomrule
\end{tabular}
\end{center}

The two lowest-correlation SEDs (exp47, exp50) share the Boredom recipe;
the addition of 2025 BG in exp50 did not further lower residual-corr but
\emph{shifted the correlation structure on hidden test rows} by enforcing
acoustic-site invariance at training time. This site-invariance is what
the LB rewards; the same-site evaluation cannot distinguish exp47 from
exp50 on residual-corr alone, but LB transfer behaviour separates them.

\subsection{Step 2 mechanism (v24 → v26): finding the weight saddle on a site-invariant SED}

exp50's contribution to the blend is not monotonic in weight. Sweep on the
$0.x P + 0.y\,\mathrm{exp50}$ family with V9 gate + Gauss:

\begin{center}
\begin{tabular}{cccc}
\toprule
Config & $W_{\mathrm{Perch}}$ & $W_{\mathrm{exp50}}$ & LB \\
\midrule
v23 & 0.8 & 0.10 & 0.928 \\
v24 & 0.8 & 0.20 & 0.930 \\
\textbf{v26} & \textbf{0.7} & \textbf{0.30} & \textbf{0.931} \\
v25 & 0.5 & 0.50 & 0.928 \\
\bottomrule
\end{tabular}
\end{center}

Concave curve with peak at v26. Below 0.30 we underuse exp50's
site-invariant rare-class signal; above it we suppress Perch's
xeno-canto-derived multi-region signal on common Aves. The 0.7--0.8
plateau on $W_{\mathrm{Perch}}$ is the load-bearing structure. v36
(2026-04-26) confirmed this from the other side: reducing
$W_{\mathrm{Perch}}$ to 0.5 to give weight to four highly-correlated SED
variants regressed −0.017 LB, because the four SEDs collectively duplicate
exp50's site-invariance contribution rather than adding new sources.

\subsection{Step 3 mechanism (v26 → v33): universal physics decoupled from site}

v33 adds, after the V9 gate but before Gauss, a per-class file-level
temporal coherence blend:
\[
\mathrm{final}[i, c] = (1 - \alpha)\,\mathrm{prob}[i, c] + \alpha \cdot \max_{j \in \mathrm{file}(i)} \mathrm{prob}[j, c]
\]
with $\alpha = 0.10$. The transformation depends only on test predictions
in the same file: no labeled SS lookup, no per-class rule, no row pattern,
no site identifier. It encodes the universal physical prior that species
presence is temporally extended within a single recording.

Local on 122 eval rows: macro $\Delta = +0.0034$, per-row Spearman vs
v26 = 0.994, Aves $\Delta = +0.007$. LB: +0.001. The $+0.003$ macro local
roughly translated $1/3$ to LB, consistent with the $1/3$--$1/4$ ratio
observed for the other two steps (e.g., v24 local macro $\Delta \sim
+0.13$ on the 11-file fair eval, LB +0.001).

\subsection{Why every other intervention failed}

Cross-referenced against the $\geq 12$ submissions that regressed (v17,
v18, v19, v22, v25, v28, v29, v32, v34, v36 plus the v6--v15 baseline
sweep): every regression matches one of three patterns:

\begin{enumerate}
  \item \textbf{Duplicated existing site-conflated signal.} v17 swapped
    SED29 for SED41f (Perch-pseudo distilled, residual-corr with Perch 0.69
    vs SED29's 0.79); the higher correlation with Perch means it adds
    redundant rather than independent information. v36 added four
    site-conflated SEDs at $W_{\mathrm{Perch}}$ = 0.5; collectively they
    duplicate exp50's site-invariance contribution rather than adding new
    sources.

  \item \textbf{Train-SS-fitted structure.} v18, v19 (27-head overlay/
    synth-aug), v22 (site prior + cluster lookup table), v28
    (27-head additive even with 2025 BG training), v29 (per-class OOF
    routing on labeled SS), v34 (mel-iVAE centroids built from labeled
    SS positives). All ten species in our 27-head are Gini=0 single-site
    (S08/S15/S19/S23 only), as confirmed by exp80a's S19-holdout test
    that dropped Insecta detection AUC to 0.073. Any structure fitted on
    these labels learns the site fingerprint.

  \item \textbf{Reduced existing site-invariance contributions.} v25 and
    v36 both reduced $W_{\mathrm{Perch}}$ to 0.5. Perch is the only
    backbone trained on multi-region xeno-canto data and is therefore the
    single largest source of site-invariance in our pipeline; reducing
    its weight directly removes this invariance.
\end{enumerate}

The successful pattern (v24, v26, v33) is the negation of all three: each
adds a new, independent, non-train-SS-fitted source of site-invariance.

\subsection{What this implies for further levers}

Searching for a fourth +0.001 step now requires identifying a fourth
source of site-invariance that is independent of:
\begin{itemize}
  \item Perch's xeno-canto multi-region pretraining
  \item exp50's 2025 BG mixup
  \item file-level temporal coherence
\end{itemize}

\noindent Concrete candidates we have not yet tested:
\begin{enumerate}
  \item \textbf{Multi-region BG mixing} (2024 + 2025 + 2023 SS pools, not
    just 2025 Colombia): broadens the geographic distribution exp50 must
    be invariant to, possibly capturing more acoustic environments.
  \item \textbf{Hour-of-day invariance via temporal augmentation} during
    SED training: hour distribution differs between train and hidden
    test, and our V9 gate's hour input is the only model component that
    sees hour information at inference.
  \item \textbf{Domain-adversarial training with site as adversary}: a
    gradient-reversal layer that explicitly removes site-discriminative
    features from the SED's representation.
  \item \textbf{Cross-window window-shift TTA} at inference time:
    averaging predictions over $\pm 1$-second shifts averages out timing
    artifacts that may be site-correlated. Has been ruled out before due
    to 90-min budget; current pipeline (v37 dev wall 511 s) suggests we
    may have headroom for 2-shift TTA on Perch alone.
\end{enumerate}

\noindent Each of these would test a specific site-invariance mechanism
distinct from the three we already have. The framework now has
falsifiable predictions: any new lever that does not fit one of these
mechanism categories should be expected to regress, regardless of local
metrics.


\section{Perch entropy + latent geometry on UNK vs KNOWN species (exp89-90)}

Motivated by why every specialized-head approach (exp44c/g, exp51, v22,
v28, v34, v36) regressed on LB, we examined Perch's internal behaviour
on rows containing species inside vs outside its 14,795-class vocabulary.

\subsection{Setup}
Of our 234 BirdCLEF 2026 species, 36 produce constant-zero logits in
Perch's output (var $< 10^{-6}$ across 739 labeled rows): 25 Insecta
sonotypes, 4 Amphibia not in iNat-bird-cluster, 2 Mammalia, 1 Reptilia,
plus 4 Aves (\texttt{grekis}, \texttt{limpki}, \texttt{74580},
\texttt{760266}) -- consistent with exp43r's earlier count of 31 unmapped
species. Row partition:
\begin{itemize}
  \item KNOWN\_ONLY (386 rows): GT contains $\geq 1$ mapped species, no unmapped.
  \item UNK\_ONLY (76 rows): only unmapped species in GT -- Perch fundamentally blind.
  \item UNK\_MIXED (277 rows): both mapped and unmapped species present.
\end{itemize}

\subsection{Finding 1: latent magnitude reduces on UNK rows}

\begin{center}
\begin{tabular}{lccc}
\toprule
Metric                & KNOWN\_ONLY (386) & UNK\_ONLY (76) & UNK\_MIXED (277) \\
\midrule
$\|\mathrm{emb}\|_2$  & $4.649 \pm 0.27$ & $4.348 \pm 0.25$ & $4.578 \pm 0.30$ \\
$\max|\mathrm{emb}_i|$ & $0.642 \pm 0.09$ & $0.591 \pm 0.09$ & $0.613 \pm 0.09$ \\
$\mathrm{Var}(\mathrm{emb})$ & $0.014 \pm 0.002$ & $0.012 \pm 0.001$ & $0.014 \pm 0.002$ \\
$\mathrm{dist}_\mathrm{ratio}$ (UNK/HIT) & $1.175 \pm 0.17$ & $0.864 \pm 0.05$ & $0.888 \pm 0.11$ \\
\bottomrule
\end{tabular}
\end{center}

UNK\_ONLY rows have measurably reduced embedding magnitude (L2 $-6.5\%$),
max-component ($-7.9\%$), and variance ($-14\%$) relative to KNOWN\_ONLY.
Centroid-distance ratio (computed with train-split-only centroids,
no leak) cleanly separates: UNK rows are closer to UNK\_ centroid,
KNOWN rows closer to HIT\_centroid (AUC $0.976$ for cross-site
separation).

\subsection{Finding 2: output is overconfident on UNK rows (output-level Anosognosia)}

\begin{center}
\begin{tabular}{lcc}
\toprule
Metric                          & KNOWN\_ONLY & UNK\_ONLY \\
\midrule
$H_{234}$ (Shannon, normalized) & $0.963 \pm 0.002$ & $\mathbf{0.960 \pm 0.001}$ \\
top1\_prob                      & $0.995 \pm 0.005$ & $\mathbf{0.997 \pm 0.004}$ \\
top5\_mass                      & $0.045 \pm 0.002$ & $\mathbf{0.047 \pm 0.002}$ \\
\bottomrule
\end{tabular}
\end{center}

Counter-intuitively, Perch is \emph{more} confident on UNK rows --
lower entropy, higher top-1 probability, more mass concentrated in the
top-5 species. The latent recognizes the input is unfamiliar
(reduced L2), but the output head suppresses this and confidently picks
some mapped species that is acoustically nearest to the unknown vocalization.

This is consistent with Perch's training objective: its 14,795-class head
was trained to assign mass to one of those classes, with no
``out-of-distribution'' option. The model lacks the language to express
``I don't know'', so it commits to a confident wrong answer. The latent
representation, which is upstream of the species head and trained
end-to-end, retains a measurable trace of the unfamiliarity.

\subsection{Finding 3: site confound dominates the centroid-distance signal}

The cross-site UNK-detection AUC of $0.976$ for embedding-distance ratio
is misleading. Site distribution of UNK\_ONLY vs KNOWN\_ONLY rows:
\begin{center}
\begin{tabular}{lcc}
\toprule
Site & UNK\_ONLY count & KNOWN\_ONLY count \\
\midrule
S22  & 0   & 344 \\
S08  & 43  & 0   \\
S09  & 14  & 4   \\
S13  & 6   & 12  \\
S18  & 8   & 4   \\
S19  & 2   & 12  \\
others & 3 & 10  \\
\bottomrule
\end{tabular}
\end{center}

S22 (344 KNOWN, 0 UNK) and S08 (43 UNK, 0 KNOWN) account for $> 80\%$ of
the comparison. The "OOD detection" AUC is largely "site classification".

\subsection{Finding 4: within-site detection is direction-inconsistent}

Restricting to sites that have both classes ($n \geq 3$ each) -- S09,
S13, S18 -- the combined uncertainty score's within-site UNK-vs-KNOWN AUC:
\begin{center}
\begin{tabular}{lccc}
\toprule
Site & $n_{\mathrm{UNK}}$ & $n_{\mathrm{KNOWN}}$ & $U_\mathrm{combined}$ AUC \\
\midrule
S09 & 14 & 4  & $\mathbf{0.339}$ (\textbf{reversed}) \\
S13 & 6  & 12 & $1.000$ \\
S18 & 8  & 4  & $0.969$ \\
\bottomrule
\end{tabular}
\end{center}

S13 and S18 show the expected direction (UNK rows have lower $\|\mathrm{emb}\|$,
higher top1\_prob), but at S09 the direction reverses: UNK rows look
\emph{less uncertain} than KNOWN rows. Mechanistic guess: at S09 (nocturnal
Insecta-heavy site), UNK\_ONLY rows have a single dominant Insecta call
that Perch confidently maps to a generic bird species, while KNOWN\_ONLY
rows happen to contain rare Aves species the model is uncertain about.

A naive routing variable derived from $U_\mathrm{combined}$ flips on a
per-site basis -- not a universal signal.

\subsection{Finding 5: routing on uncertainty regresses}

Tested four routing schemes on 122 eval rows: hard threshold ($U > 0.5,
1.0, 1.5$), sigmoid-gated continuous re-weighting (scale 0.5--3.0), and
target-routing to exp84b (external-fine-tuned SED). Every variant
regressed:

\begin{center}
\begin{tabular}{lcccccc}
\toprule
Config                      & macro $\Delta$ & sp\_row & Aves $\Delta$ & Insecta $\Delta$ & Mammalia $\Delta$ & Reptilia $\Delta$ \\
\midrule
v33 ref                     & 0      & 1.000 & 0     & 0      & 0     & 0 \\
Hard $U > 1.5$, n=67       & $-0.029$ & 0.999 & $+0.003$ & $-0.072$ & $+0.008$ & $-0.136$ \\
$\sigma$-route scale=1.0    & $-0.052$ & 0.999 & $+0.015$ & $-0.112$ & $-0.036$ & $-0.330$ \\
$\sigma$-route → exp84b s=2 & $-0.065$ & 0.999 & $+0.022$ & $-0.147$ & $-0.059$ & $-0.324$ \\
\bottomrule
\end{tabular}
\end{center}

\paragraph{Mechanism of failure.}
The diagnostic is correct (Perch latent shows reduced magnitude on UNK rows)
but there is no \emph{trustworthy alternative model} to route uncertain rows
to. Every SED we have (exp50, exp59, exp73, exp84b, exp83) was trained on
the same 4-Insecta-site labeled SS data. Routing uncertain rows to any of
them amplifies the same site shortcut that hurt v18, v19, v22, v28, v34, v36.
The diagnostic is real but the treatment is not available within our
current model zoo.

\subsection{What this changes for the search}

The framework predicts a fourth source of site-invariance is needed; this
section quantifies the missing component as an "invariant non-Aves
predictor". Until such a predictor exists, MoE-style routing on Perch
uncertainty is structurally blocked.

A subtle further hypothesis to test: even though absolute uncertainty
direction reverses across sites (S09 vs S13/S18), the \emph{deviation
from per-site mean confidence} may be consistently signed -- i.e., UNK
rows may always be in the upper tail of within-site confidence, even
when absolute confidence varies by site. exp91 will run this normalised
comparison; if signal persists, the routing variable is per-site
z-scored confidence, not raw confidence.


\subsection{exp91: within-site z-scoring partially recovers the signal}

User hypothesis after seeing exp89: even though absolute confidence
direction reverses across sites (S09 reverse), the per-site z-scored
deviation may be consistently signed -- UNK rows in the upper tail of
each site's confidence distribution.

Tested by per-site z-scoring 8 metrics (entropy, top1, top5\_mass,
top1−top2 margin, embedding L2, max, var, participation ratio) on the
739 labeled rows.

\paragraph{Population-level signal exists but is small.}
$U_\mathrm{combined}$ (combined per-site z-scored uncertainty)
class means:

\begin{center}
\begin{tabular}{lc}
\toprule
Class      & $U_\mathrm{combined}$ mean \\
\midrule
UNK\_ONLY   & $+0.181$ \\
UNK\_MIXED  & $+0.029$ \\
KNOWN\_ONLY & $-0.056$ \\
\bottomrule
\end{tabular}
\end{center}

Clean monotonic ordering UNK\_ONLY $>$ UNK\_MIXED $>$ KNOWN\_ONLY:
overconfidence-on-UNK is a systematic signal at population level. The
direction inconsistency at S09 is a small-sample (n=4 KNOWN) artifact
within an underlying population-level effect.

Pooled (within qualifying sites) UNK\_ONLY-vs-KNOWN\_ONLY z-scored AUCs:
\begin{itemize}
  \item $\mathrm{emb}\_\mathrm{max}$: AUC $0.275$ (strength $0.45$, inverse direction)
  \item $\mathrm{emb}\_\mathrm{var}$: AUC $0.334$ (strength $0.33$, inverse)
  \item $\mathrm{emb}_2$: AUC $0.332$ (strength $0.34$, inverse)
  \item participation ratio: AUC $0.612$ (strength $0.22$)
  \item top1: AUC $0.595$ (strength $0.19$)
\end{itemize}

\paragraph{Routing tests.}
With per-site z-scored uncertainty $U_z$ as the routing variable:

\begin{center}
\begin{tabular}{lcccc}
\toprule
Variant                                  & macro $\Delta$ & sp\_row & Aves $\Delta$ & predicted LB \\
\midrule
v33 ref                                  & 0       & 1.000  & 0       & ref \\
\textbf{C-suppress(map) $U_z > 1.0$ scale=0.3} & $\mathbf{+0.00070}$ & $\mathbf{0.99992}$ & $\mathbf{+0.0019}$ & A \\
C-suppress scale=0.2                     & $+0.00056$ & 0.99995 & $+0.0015$ & A \\
C-suppress scale=0.1                     & $+0.00040$ & 0.99999 & $+0.0016$ & A \\
A-damp(global) scale=0.05                & $-0.029$ & 1.000 & $+0.005$ & C \\
B-route-to-exp50 $U_z > 0.5$ (n=123)     & $-0.050$ & 0.981 & $-0.046$ & C \\
\bottomrule
\end{tabular}
\end{center}

\paragraph{First class-A positive candidate since v33.}
``C-suppress'' soft-suppresses mapped-species columns on uncertain rows
($U_z > 1.0$): $\mathrm{prob}[i, c] \mathrel{*}= (1 - \mathrm{scale})$
for $c \in \mathrm{mapped\_idx}$. Mechanism: it counters Perch's
``overconfident-wrong on UNK'' specifically by dampening only the
mapped-species channels that Perch is overconfident about, leaving the
other 36 unmapped channels and the underlying file-max coherence
intact.

The local effect is small ($\sim 1/5$ of v33's local gain). LB-transfer
estimate using v33's gain ratio: $+0.0002$ to $+0.0003$, near noise floor.

\paragraph{Variant A and B failure modes.}
A-damp regresses by $-0.029$ macro because flat damping reduces
ranking signal globally on uncertain rows. B-route to exp50 regresses
by $-0.05$ because routing UNK rows to a SED with site-conflated
labels amplifies the same site shortcut documented in v22, v28, v34.

\subsection{What this leaves open}

Two directions remain unexplored:
\begin{itemize}
  \item \textbf{Larger-scale C-suppress sweep.} Whether $U_z$ threshold
    or scale near $1.0$ produces measurable improvement without
    breaking sp\_row.
  \item \textbf{Spectral (FFT) features of mel-spectrogram time axis.}
    The pooled mel cache is $T = 16$ frames; FFT along time provides
    $\sim 8$ frequency bins describing temporal modulation patterns.
    Insecta sonotypes are sustained narrow-band; bird calls are
    transient broadband. Different temporal modulation signatures may
    enable UNK detection independent of site.
\end{itemize}


\section{exp95--96: per-(row, class) selective lever from teacher-vote signals (2026-04-27)}

After exp89--94 confirmed that Perch's latent has TN-vs-FN separability
(but SED50's does not, and ConvNeXt-tiny's does), exp95 ran an exhaustive
signal hunt over 16 candidate features for distinguishing the four
TP/FP/TN/FN quadrants on (row, class) pairs.

\subsection{exp95: signal ranking}

Per-feature AUC for separating each pair, computed over 215 TP / 638 FP
/ 10061 TN / 171 FN pairs across 15 mapped Aves classes:

\begin{center}
\begin{tabular}{lccc}
\toprule
Feature                               & TP\_vs\_FP & \textbf{TN\_vs\_FN} & CORR\_vs\_WRONG \\
\midrule
\textbf{sed\_on\_c (exp50 raw on class)}     & 0.053 & \textbf{0.983} & 0.632 \\
\textbf{cnxt\_on\_c (exp59 raw on class)}    & 0.050 & \textbf{0.982} & 0.676 \\
v33\_on\_c (self-prediction)                  & 0.244 & 0.925 & 0.968 \\
\textbf{perch\_sed\_disagreement}             & \textbf{0.890} & 0.624 & 0.850 \\
perch\_on\_c                                  & 0.666 & 0.255 & 0.786 \\
file\_max\_perch\_c                           & 0.558 & 0.268 & 0.758 \\
perch top1                                    & 0.437 & 0.740 & 0.612 \\
Perch spatial-FFT ac\_dc\_ratio               & 0.541 & 0.685 & 0.682 \\
Perch spatial-FFT low\_high\_ratio            & 0.506 & 0.699 & 0.648 \\
ConvNeXt pre-head FFT ac\_dc\_ratio           & 0.522 & 0.673 & 0.658 \\
SED50 pre-head FFT ac\_dc\_ratio              & 0.542 & 0.354 & 0.453 \\
\bottomrule
\end{tabular}
\end{center}

\paragraph{Two near-perfect signals.}
\begin{itemize}
  \item \texttt{sed\_on\_c} (exp50's raw prediction on class $c$ at row
    $i$): TN\_vs\_FN AUC $= 0.983$. When v33 says low on class $c$ but
    exp50 says high, the truth is almost certainly that $c$ is positive
    (FN, not TN). Mean $\mathrm{exp50}_\mathrm{TN} = 0.005$ vs
    $\mathrm{exp50}_\mathrm{FN} = 0.65$ — 130$\times$ separation.
  \item \texttt{perch\_sed\_disagreement} ($|\mathrm{Perch}[i,c] -
    \mathrm{exp50}[i,c]|$): TP\_vs\_FP AUC $= 0.890$. When v33 says high
    and the two teachers disagree, the prediction is likely FP.
    Mean disagreement $\mathrm{TP} = 0.36$ vs $\mathrm{FP} = 0.82$.
\end{itemize}

\paragraph{The transformer-vs-CNN hypothesis is partially refuted.}
ConvNeXt-tiny (pure CNN, ImageNet-22k pretrained) shows nearly identical
latent FN signature as Perch (ac\_dc\_ratio TN\_vs\_FN AUC 0.673 vs
Perch's 0.685). SED50 (pure CNN, HGNet-B0, 4.9M params, BatchNorm) is
the outlier with flat AUC ($\sim 0.35$). The distinguishing factor is
not transformer-vs-CNN but rather \emph{rich pretraining + modern
normalization (LayerNorm) + sufficient depth}.

\subsection{exp96: per-(row, class) selective lever}

Two surgical operations on v33, both per-(row, class) — neither modifies
$W_\mathrm{Perch}$ globally (preserving the v36 invariant):

\paragraph{FN\_RESCUE.} Where $\mathrm{v33}[i, c] < \tau_\mathrm{low}$
AND $\max(\mathrm{exp50}, \mathrm{exp59})[i, c] > \delta_t$:
\[
\mathrm{new}[i, c] = \mathrm{v33}[i, c] + \alpha_\mathrm{fn} \cdot
\mathrm{teacher}[i, c] \cdot (1 - \mathrm{v33}[i, c])
\]

\paragraph{FP\_SUPPRESS.} Where $\mathrm{v33}[i, c] > \tau_\mathrm{high}$
AND $|\mathrm{Perch}[i, c] - \mathrm{exp50}[i, c]| > \delta_d$:
\[
\mathrm{new}[i, c] = \mathrm{v33}[i, c] \cdot (1 - \alpha_\mathrm{fp}
\cdot \mathrm{disagree}_{\mathrm{normalized}})
\]

\paragraph{Sweep results on 122 eval rows.}
\begin{center}
\begin{tabular}{lcccccc}
\toprule
Config                                    & macro $\Delta$ & sp\_row & Aves $\Delta$ & Insecta $\Delta$ & Mammal $\Delta$ & Amphibia $\Delta$ \\
\midrule
v33 ref                                   & 0       & 1.000 & 0       & 0       & 0       & 0 \\
\textbf{FN\_R $\tau_l = 0.4$, $\alpha = 0.5$, $\delta = 0.3$} & $\mathbf{+0.0068}$ & $\mathbf{0.999}$ & $\mathbf{+0.009}$ & $\mathbf{+0.008}$ & $+0.004$ & $+0.005$ \\
FN\_R $\tau_l = 0.5$, $\alpha = 0.5$       & $+0.0065$ & 0.999 & $+0.009$ & $+0.008$ & $+0.004$ & $+0.005$ \\
FN+FP $\tau_l = 0.4, \alpha_\mathrm{fn} = 0.3, \tau_h = 0.5, \alpha_\mathrm{fp} = 0.2$ & $+0.0065$ & 0.999 & $+0.007$ & $+0.006$ & $+0.004$ & $+0.007$ \\
FN\_R(exp50 only) $\tau_l = 0.5$, $\alpha = 0.3$ & $+0.0033$ & 1.000 & $0$    & $+0.006$ & $+0.004$ & $+0.003$ \\
\bottomrule
\end{tabular}
\end{center}

\paragraph{Comparison with prior submissions.}
\begin{center}
\begin{tabular}{lccccc}
\toprule
                  & macro\_d local & sp\_row & Aves $\Delta$ & $W_\mathrm{Perch}$ & LB outcome \\
\midrule
v33 (positive)    & $+0.003$       & 0.994   & $+0.007$ & 0.7 & $+0.001$ \\
v34 (negative)    & $+0.026$       & 0.990   & $+0.008$ & 0.7 (+ mel-iVAE) & $-0.016$ \\
v36 (negative)    & $+0.013$       & 0.999   & $+0.024$ & \textbf{0.5 (violated)} & $-0.017$ \\
\textbf{Best exp96 (untested)} & \textbf{+0.007} & \textbf{0.999} & \textbf{+0.009} & \textbf{0.7 (preserved)} & ? \\
\bottomrule
\end{tabular}
\end{center}

The candidate's profile most closely matches v33's (the only LB-positive
modification): roughly twice the local macro gain, identical sp\_row
safety margin, comparable Aves $\Delta$. Crucially preserves
$W_\mathrm{Perch} \geq 0.7$ globally.

\paragraph{Why this candidate is structurally different from v34/v36.}
\begin{itemize}
  \item v34: introduced mel-iVAE as a new universal lever applied to all
    234 columns globally. Inherits site shortcut from train-SS-fitted
    centroids.
  \item v36: rebalanced the entire blend (4 SEDs at 0.125 each,
    $W_\mathrm{Perch} = 0.5$). Amplified site shortcut by overweighting
    site-conflated SEDs.
  \item exp96 best: surgical per-(row, class) modification only where
    multiple teacher signals provide additional information. Surface
    area is small (only rows where v33 is uncertain on $c$ and exp50/exp59
    have positive signal). $W_\mathrm{Perch}$ stays at 0.7. Local profile
    twice v33's gain with same safety margin.
\end{itemize}

\paragraph{Residual transfer risk.}
The mechanism uses \texttt{sed\_on\_c} as an FN signal. exp50 and exp59
are both train-SS-trained and inherit the 4-Insecta-site shortcut.
Therefore on hidden test sites, an "exp50 says high but v33 says low"
event may reflect site fingerprint rather than missed species. The
surgical surface area limits this risk but does not eliminate it.

The framework predicts class-A transfer (v33-like profile, preserved
$W_\mathrm{Perch}$, universal mechanism, all-taxa $\geq 0$ Aves $> 0$),
but the $-0.017$ catastrophe of v36 with similar formal class-A profile
is a warning. LB testing is the only definitive validation.


\subsection{exp97: robustness checks on the best lever}

Five strain tests on the FN\_RESCUE (best from exp96):

\paragraph{R1. Train-vs-eval consistency.} The lever produces $\Delta$
macro $= +0.0068$ on 122 eval rows (40 evaluable classes) but only
$+0.0007$ on 617 train rows (66 classes, mostly saturated at AUC
$\sim 0.98$). Train/eval gain ratio = $0.11$ — the lever is not fitting
to train-distribution patterns. Low ratio is consistent with the gain
coming from rare-class rescue rather than general-class tuning.

\paragraph{R2. Per-site analysis.} Eval $\Delta$ macro by site:

\begin{center}
\begin{tabular}{lcccc}
\toprule
Site (eval) & n\_rows & v33 & lever & $\Delta$ \\
\midrule
S08 (Insecta-dominant) & 12 & 0.619 & 0.596 & $\mathbf{-0.023}$ \\
S13                    & 12 & 0.972 & 0.972 & 0.000 \\
S19                    & 12 & 0.535 & 0.535 & 0.000 \\
S22 (Aves-dominant)    & 60 & 0.849 & 0.854 & $+0.005$ \\
S23                    & 24 & 0.556 & 0.594 & $\mathbf{+0.039}$ \\
\bottomrule
\end{tabular}
\end{center}

The +0.0068 net gain comes mostly from S23 ($+0.039$, 24 rows). S08
($-0.023$, 12 rows) and S22 ($+0.005$, 60 rows) push back. S13/S19 stay
flat. The S08 regression is a warning: the lever boosts mapped-Aves
predictions on rows that have only unmapped Insecta GT, which can hurt
some Aves classes' AUC by adding incorrect-positive rankings.

\paragraph{R3. Per-class attribution.} Top contributors to eval gain:
22967 Amphibia ($+0.091$), wfwduc1 Aves ($+0.074$), 22973 Amphibia
($+0.054$), 47158son11 Insecta ($+0.042$). Top regressors: 555146
Amphibia ($-0.043$, already saturated at v33 = 0.985), 517063 Amphibia
($-0.025$, already at 0.993). The lever helps weak classes but slightly
hurts saturated ones — net positive but with cost.

\paragraph{R4. Holdout-site sensitivity.} Drop one site at a time, recompute:

\begin{center}
\begin{tabular}{lc}
\toprule
Holdout & $\Delta$ macro on remaining \\
\midrule
S08 & $+0.007$ \\
S09 & $+0.007$ \\
S13 & $+0.008$ \\
S19 & $+0.005$ \\
S22 & $+0.013$ \\
S23 & $+0.004$ (smallest — confirms S23 leads contribution) \\
\bottomrule
\end{tabular}
\end{center}

All holdouts retain positive net gain. Lever is not a single-site
phenomenon. With S22 held out the gain grows to $+0.013$, indicating
S22's "common Aves" rows actually mute the gain compared to rare-class
sites.

\paragraph{R5. Finer parameter sweep.} A wider sweep around the
exp96 best identified an even better config: $\tau_l = 0.55$,
$\alpha = 0.7$, $\delta = 0.2$ — macro $\Delta = +0.010$, sp\_row 0.997,
Aves $\Delta = +0.017$, all five taxa $\geq 0$.

\begin{center}
\begin{tabular}{lcccccc}
\toprule
                                       & macro $\Delta$ & sp\_row & Aves $\Delta$ & W\_Perch & ratio (Aves $\Delta$ / v33) & predicted \\
\midrule
v33 (LB +0.001 reference)              & $+0.003$ & 0.994   & $+0.007$ & 0.7 & 1.0 & ref \\
v34 (LB $-0.016$)                       & $+0.026$ & 0.990   & $+0.008$ & 0.7 (+ mel-iVAE) & 1.1 & X \\
v36 (LB $-0.017$)                       & $+0.013$ & 0.999   & $+0.024$ & \textbf{0.5} & 3.4 & X \\
\textbf{exp97 best (untested)}        & $+0.010$ & 0.997   & $+0.017$ & 0.7 (preserved) & 2.4 & ? \\
exp96 best ($\tau_l=0.4, \alpha=0.5, \delta=0.3$) & $+0.007$ & 0.999 & $+0.009$ & 0.7 & 1.3 & ? \\
\bottomrule
\end{tabular}
\end{center}

The exp97 best lies between v33 and v36 in Aves $\Delta$ magnitude
($\times 2.4$ v33's). The exp96 best is closer to v33 ($\times 1.3$).
Both preserve $W_\mathrm{Perch}$. Neither has been LB-tested.

\subsection{Decision space}

\paragraph{Bull case.}
\begin{itemize}
  \item Mechanism is genuinely new (per-(row, class) teacher-vote rescue).
  \item Holdout-site robust: not driven by any single site.
  \item Train/eval ratio 0.11 — not train-distribution-fitted.
  \item W\_Perch preserved at 0.7.
  \item v33 transferred at $+0.001$ with similar profile, this is twice
    the local magnitude.
  \item Best-case LB transfer: $+0.002$ to $+0.004$.
\end{itemize}

\paragraph{Bear case.}
\begin{itemize}
  \item S08 regresses by $-0.023$. The Insecta-only site is exactly where
    Perch's overconfident-wrong pattern is strongest, and the lever
    sometimes amplifies that pattern.
  \item exp50 and exp59 (the teachers we trust) both inherit the 4-Insecta-
    site shortcut. Their per-class predictions on hidden test sites may
    be site-fingerprint signals rather than species signals.
  \item v36 had a similar formal profile (sp\_row 0.999, all taxa $\geq 0$,
    Aves $+0.024$) and regressed by $-0.017$. The structural similarity
    is concerning.
  \item Worst-case LB transfer: $-0.005$ to $-0.015$.
\end{itemize}

\paragraph{Decision.}
With 4 LB slots remaining today and high uncertainty, the candidate is
recorded but not yet submitted. The two-sided risk distribution is
honest acknowledgment that local metrics have repeatedly failed to
predict LB transfer for this lever family.


\section{exp99--100: LR-based correction lever — strongest cross-site validation yet}

After exp93--98 identified per-class teacher signals (\texttt{sed\_on\_c},
\texttt{cnxt\_on\_c}, \texttt{perch\_sed\_disagree}) with near-perfect
TP/FP/TN/FN separability, exp99 fitted multivariate logistic regression
detectors and exp100 deployed them as a per-(row, class) correction lever.

\subsection{exp99: LR FP/FN detectors}

For each (row, class) pair where v33 fires above 0.5 (FP territory) or
below 0.5 (FN territory), built features:
\begin{itemize}
  \item Universal: \texttt{perch\_on\_c}, \texttt{exp50\_on\_c},
    \texttt{exp59\_on\_c}, \texttt{perch\_sed\_disagree},
    \texttt{file\_mean\_e50\_c}, \texttt{file\_std\_e50\_c},
    \texttt{file\_uniform\_e50\_c}, \texttt{v33\_on\_c}
  \item Site-specific: \texttt{class\_dom\_site\_match},
    \texttt{class\_site\_concentration}, \texttt{hour\_match\_strength}
\end{itemize}

Trained LR on TRAIN pairs (673 FP, 8582 FN), evaluated on EVAL pairs
(180 FP, 1650 FN):

\begin{center}
\begin{tabular}{lcccc}
\toprule
Detector & Features & Train AUC & Eval AUC \\
\midrule
FP & FULL (incl. site)         & 0.979 & 0.949 \\
FP & \textbf{UNIVERSAL only}   & 0.980 & \textbf{0.909} \\
FN & FULL                       & 0.998 & 0.912 \\
FN & \textbf{UNIVERSAL only}   & 0.998 & \textbf{0.987} \\
\bottomrule
\end{tabular}
\end{center}

Universal-only features retain almost all signal. Site-specific features
add little (and FN detector improves with universal-only — site features
were overfitting).

\paragraph{Key driver: file-level uniformity.} The strongest "↑FP"
coefficient is \texttt{file\_uniform\_e50\_c} (+0.56): when exp50 fires
uniformly across all 12 windows of a file, that's a site fingerprint
(species calls are sparse and bursty); when it fires on just a few
windows, that's a real species detection.

\subsection{exp100: LOSO-site cross-validation — decisive evidence of generalization}

The strongest test we have ever conducted of cross-site generalization.
Train LR on labeled-SS sites except one, evaluate on the held-out site:

\begin{center}
\begin{tabular}{lccc}
\toprule
Holdout site & $n$ pairs & FP detector AUC & FN detector AUC \\
\midrule
S08 & 55 / 845    & 1.000  & 0.999 \\
S09 & 22          & 0.950  & --    \\
S15 & 143 / 577   & 0.966  & 0.986 \\
S19 & --          & --     & 0.998 \\
S22 & 525 / 6630  & 0.952  & 0.996 \\
S23 & 91 / 449    & 0.999  & 0.929 \\
\midrule
Mean LOSO         &        & \textbf{0.973} & \textbf{0.981} \\
\bottomrule
\end{tabular}
\end{center}

This is the strongest cross-site evidence we have collected. Across all
5 (FP) / 5 (FN) holdout configurations, the detectors retain AUC $\geq
0.93$. The LR is genuinely using site-invariant signals (file uniformity,
model agreement, file mean) -- not memorizing site-class associations.

\subsection{Application to v33}

Sweep $\alpha_\mathrm{fn}$ (FN boost) $\times$ $\beta_\mathrm{fp}$ (FP suppress)
on 122 eval rows:

\begin{center}
\begin{tabular}{lcccccccc}
\toprule
Config & macro $\Delta$ & sp\_row & Aves $\Delta$ & Insecta & Mamm & Amph & Rept & predicted \\
\midrule
v33 ref & 0     & 1.000  & 0       & 0 & 0 & 0 & 0 & ref \\
\textbf{LR $\alpha=0.5$ $\beta=0.1$} & $\mathbf{+0.005}$ & $\mathbf{0.999}$ & $\mathbf{+0.022}$ & 0 & 0 & 0 & 0 & A \\
LR $\alpha=0.5$ $\beta=0.0$           & $+0.005$ & 0.999  & $+0.022$ & 0 & 0 & 0 & 0 & A \\
LR $\alpha=0.3$ $\beta=0.1$           & $+0.0035$ & 0.9996 & $+0.016$ & 0 & 0 & 0 & 0 & A \\
LR $\alpha=0.2$ $\beta=0.0$           & $+0.0024$ & 0.9999 & $+0.011$ & 0 & 0 & 0 & 0 & A \\
\bottomrule
\end{tabular}
\end{center}

Best: $\alpha = 0.5$, $\beta = 0.1$. Aves $\Delta = +0.022$ is the entire
gain (other taxa $= 0$). Per-site analysis shows S08 $\Delta = 0$
(no spurious modification), S23 $\Delta = +0.028$ (the gain
concentrates on the rare-Aves-rich site).

\subsection{Honest LB-transfer assessment}

\paragraph{Bull-case factors.}
\begin{itemize}
  \item LOSO-site mean AUC $= 0.973$/$0.981$ — strongest cross-site
    evidence we have ever obtained.
  \item Universal-only features retain signal ($> 0.91$ AUC); site-
    identity-dependent features are not needed.
  \item $W_\mathrm{Perch} = 0.7$ globally preserved (per-class
    adjustment).
  \item exp97's S08 regression $-0.023$ becomes 0 here (per-site
    analysis confirms).
  \item macro $\Delta = +0.005$ is half of exp97's best and a third of
    v36's. If transfer ratio holds, predicted LB $\sim +0.001$.
\end{itemize}

\paragraph{Bear-case factors.}
\begin{itemize}
  \item Aves $\Delta = +0.022$ is structurally similar to v36's $+0.024$
    that crashed at LB $-0.017$. Past patterns of high Aves $\Delta$
    transferring to negative LB are an active warning.
  \item LR was fitted on labeled SS pairs. Even with universal features,
    the LR coefficients reflect SS distribution. Hidden-test feature
    distribution may shift.
  \item Insecta/Mammalia/Amphibia/Reptilia $\Delta = 0$ — all gain
    is on Aves, the same taxon where v36 also lifted.
  \item Two consecutive predictions of "class A LB-positive" (v34, v36)
    have failed.
\end{itemize}

\paragraph{Verdict.} This is the strongest LB candidate we have produced
since v33 became production. The LOSO-site cross-validation is qualitative
evidence the previous candidates lacked. The mechanism is genuinely
different from v36 (per-class targeted using LR-fitted signals, vs
uniform global SED rebalancing). If LB framework predictions are correct,
$+0.001$ to $+0.003$. If the past anti-correlation pattern continues,
$-0.005$ to $-0.010$. Worth one of the four remaining slots, but with
clear-eyed acknowledgment of uncertainty.


\section{Final LB sweep (2026-04-27): all 5 slots regressed, framework refuted}

After the LR-correction lever achieved LOSO-site CV AUC 0.97 (the
strongest cross-site validation we had ever produced), we used all 5
daily LB submission slots to test orthogonal hypotheses about which
component of the lever family transfers.

\paragraph{Submissions.}
\begin{center}
\begin{tabular}{lllll}
\toprule
Sub & Config                                  & local macro $\Delta$ & local Aves $\Delta$ & LB \\
\midrule
v33 (production ref) & 0.7P + 0.3 exp50 + V9 + file-max $\alpha=0.10$ & ref & ref & 0.932 \\
v38 & LR $\alpha_\mathrm{fn}=0.3, \beta_\mathrm{fp}=0.1$ (conservative) & $+0.0035$ & $+0.016$ & $0.920$ \\
v39 & LR $\alpha=0.5, \beta=0.1$ (aggressive)                            & $+0.005$  & $+0.022$ & $0.915$ \\
v40 & LR $\alpha=0.5, \beta=0.0$ (FN-only ablation)                       & $+0.005$  & $+0.022$ & $0.917$ \\
v41 & exp97 threshold rule (no LR)                                       & $+0.010$  & $+0.017$ & $0.918$ \\
v42 & file-max $\alpha=0.20$ (universal physics scaled)                  & $+0.006$  & $+0.014$ & $0.918$ \\
\bottomrule
\end{tabular}
\end{center}

\paragraph{Results: all 5 regressed.}
LB outcomes range $-0.012$ to $-0.017$. The LR-correction lever
(v38--v40), the simpler threshold-rule equivalent (v41), and even the
universal-physics scaling (v42) all crashed by similar magnitude.

\paragraph{LOSO-site CV invalidated.} The strongest cross-site
validation we ever produced (LR detector LOSO-site mean AUC $0.973$
for FP, $0.981$ for FN) failed to predict LB transfer. This is the
clearest possible refutation of our extended anti-correlation rule:
\emph{cross-site validation on labeled SS holdouts is fundamentally
disconnected from hidden-test-site generalization}.

\paragraph{LR vs threshold-rule: equivalent failures.}
\begin{itemize}
  \item v40 (LR FN-only): LB $-0.015$
  \item v41 (exp97 threshold rule, no LR): LB $-0.014$
\end{itemize}
The LR-fitting added no measurable value beyond a simple threshold rule
applied to raw teacher scores. The problem is the FAMILY (per-(row,
class) teacher-boost on low v33), not the calibration method.

\paragraph{FP-suppress component is dead weight.}
\begin{itemize}
  \item v39 (LR with $\beta=0.1$ FP-suppress): LB $-0.017$
  \item v40 (LR same $\alpha$ but $\beta=0$): LB $-0.015$
\end{itemize}
Adding FP-suppress costs an extra $-0.002$. Local TP/FP separability
AUC $0.89$ does not translate to LB benefit; in fact it actively hurts.

\paragraph{Universal physics also non-monotonic.}
v33 uses file-max $\alpha = 0.10$ and was LB $+0.001$. v42 with
$\alpha = 0.20$ crashed to LB $-0.014$. \textbf{file-max coherence has
a sharp peak at $\alpha = 0.10$ and degrades rapidly above it.} Even
the most defensible "universal physics" lever has a narrow viable range.

\subsection{Updated lever taxonomy after final sweep}

11 distinct LB-tested mechanisms, sorted by LB $\Delta$:

\begin{center}
\begin{tabular}{lr}
\toprule
Mechanism                                           & LB $\Delta$ vs ref \\
\midrule
\textbf{Universal physics @ optimal $\alpha$} (v33 file-max 0.10) & $\mathbf{+0.001}$ \\
R5 routing (v32, model-property based)              & $-0.001$ \\
27-head additive (v28)                              & $-0.002$ \\
Per-class OOF routing (v29)                         & $-0.004$ \\
SED swap to Perch-pseudo SED (v17 SED41f)           & $-0.007$ \\
LR conservative (v38)                               & $-0.012$ \\
\textbf{Universal physics scaled} (v42 $\alpha=0.20$) & $-0.014$ \\
Threshold rule (v41)                                & $-0.014$ \\
LR FN-only (v40)                                    & $-0.015$ \\
Train-SS centroid (v34 mel-iVAE)                    & $-0.016$ \\
Train-SS lookup (v22 site prior + cluster)          & $-0.016$ \\
LR aggressive (v39)                                 & $-0.017$ \\
Global rebalance W\_Perch=0.5 (v36)                  & $-0.017$ \\
\bottomrule
\end{tabular}
\end{center}

\paragraph{Across 11 mechanisms, only v33's exact configuration is
LB-positive.} Every variation -- whether per-class fitted, threshold
rule, universal physics scaled, train-SS lookup, global rebalance, or
SED swap -- regresses by $-0.012$ to $-0.017$ from v33's 0.932. The
implication is that v33 is not just a local optimum but a true ceiling
for our current model zoo.

\paragraph{Falsifiable predictions all failed.} Our framework predicted
"class A" for v34, v36, v38, v39, v40, v41 -- six consecutive
predictions. All six regressed. The framework as currently formulated
is not predictive of LB transfer.

\subsection{What this leaves}

Three remaining options:
\begin{enumerate}
  \item \textbf{Data augmentation}: add training-time site-invariance
    via multi-region BG mixing (2023 + 2024 + 2025 SS pools), expanded
    xeno-canto / iNat downloads (especially Mammalia and Reptilia where
    we have nothing). This adds a NEW source of site-invariance
    \emph{at training time}, not at inference.
  \item \textbf{New foundation model}: Perch v3 or a transformer-based
    bioacoustic teacher with rich multi-region pretraining. Would
    occupy the $W \geq 0.7$ slot currently held by Perch v2, with
    stronger species-discriminability geometry.
  \item \textbf{Accept v33 as production ceiling.} Twenty-one LB
    submissions across two months have produced exactly one positive-
    transfer modification (v33 file-max $\alpha=0.10$). The remaining
    inference-time lever space is exhausted.
\end{enumerate}

The negative-result thesis is now complete: under our data structure
(4-Insecta-site labeled SS, multi-task SED variants all on
HGNet-B0-style architectures, single multi-region pretrained Perch
backbone), the LB-transferable lever space at inference time consists
solely of universal physical priors at narrow optimal hyperparameters.
Further LB gains require enriching the training-time data distribution
or replacing the foundation model.
